% Options for packages loaded elsewhere
% Options for packages loaded elsewhere
\PassOptionsToPackage{unicode}{hyperref}
\PassOptionsToPackage{hyphens}{url}
\PassOptionsToPackage{dvipsnames,svgnames,x11names}{xcolor}
%
\documentclass[
  letterpaper,
  DIV=11,
  numbers=noendperiod]{scrreprt}
\usepackage{xcolor}
\usepackage{amsmath,amssymb}
\setcounter{secnumdepth}{5}
\usepackage{iftex}
\ifPDFTeX
  \usepackage[T1]{fontenc}
  \usepackage[utf8]{inputenc}
  \usepackage{textcomp} % provide euro and other symbols
\else % if luatex or xetex
  \usepackage{unicode-math} % this also loads fontspec
  \defaultfontfeatures{Scale=MatchLowercase}
  \defaultfontfeatures[\rmfamily]{Ligatures=TeX,Scale=1}
\fi
\usepackage{lmodern}
\ifPDFTeX\else
  % xetex/luatex font selection
\fi
% Use upquote if available, for straight quotes in verbatim environments
\IfFileExists{upquote.sty}{\usepackage{upquote}}{}
\IfFileExists{microtype.sty}{% use microtype if available
  \usepackage[]{microtype}
  \UseMicrotypeSet[protrusion]{basicmath} % disable protrusion for tt fonts
}{}
\makeatletter
\@ifundefined{KOMAClassName}{% if non-KOMA class
  \IfFileExists{parskip.sty}{%
    \usepackage{parskip}
  }{% else
    \setlength{\parindent}{0pt}
    \setlength{\parskip}{6pt plus 2pt minus 1pt}}
}{% if KOMA class
  \KOMAoptions{parskip=half}}
\makeatother
% Make \paragraph and \subparagraph free-standing
\makeatletter
\ifx\paragraph\undefined\else
  \let\oldparagraph\paragraph
  \renewcommand{\paragraph}{
    \@ifstar
      \xxxParagraphStar
      \xxxParagraphNoStar
  }
  \newcommand{\xxxParagraphStar}[1]{\oldparagraph*{#1}\mbox{}}
  \newcommand{\xxxParagraphNoStar}[1]{\oldparagraph{#1}\mbox{}}
\fi
\ifx\subparagraph\undefined\else
  \let\oldsubparagraph\subparagraph
  \renewcommand{\subparagraph}{
    \@ifstar
      \xxxSubParagraphStar
      \xxxSubParagraphNoStar
  }
  \newcommand{\xxxSubParagraphStar}[1]{\oldsubparagraph*{#1}\mbox{}}
  \newcommand{\xxxSubParagraphNoStar}[1]{\oldsubparagraph{#1}\mbox{}}
\fi
\makeatother

\usepackage{color}
\usepackage{fancyvrb}
\newcommand{\VerbBar}{|}
\newcommand{\VERB}{\Verb[commandchars=\\\{\}]}
\DefineVerbatimEnvironment{Highlighting}{Verbatim}{commandchars=\\\{\}}
% Add ',fontsize=\small' for more characters per line
\usepackage{framed}
\definecolor{shadecolor}{RGB}{241,243,245}
\newenvironment{Shaded}{\begin{snugshade}}{\end{snugshade}}
\newcommand{\AlertTok}[1]{\textcolor[rgb]{0.68,0.00,0.00}{#1}}
\newcommand{\AnnotationTok}[1]{\textcolor[rgb]{0.37,0.37,0.37}{#1}}
\newcommand{\AttributeTok}[1]{\textcolor[rgb]{0.40,0.45,0.13}{#1}}
\newcommand{\BaseNTok}[1]{\textcolor[rgb]{0.68,0.00,0.00}{#1}}
\newcommand{\BuiltInTok}[1]{\textcolor[rgb]{0.00,0.23,0.31}{#1}}
\newcommand{\CharTok}[1]{\textcolor[rgb]{0.13,0.47,0.30}{#1}}
\newcommand{\CommentTok}[1]{\textcolor[rgb]{0.37,0.37,0.37}{#1}}
\newcommand{\CommentVarTok}[1]{\textcolor[rgb]{0.37,0.37,0.37}{\textit{#1}}}
\newcommand{\ConstantTok}[1]{\textcolor[rgb]{0.56,0.35,0.01}{#1}}
\newcommand{\ControlFlowTok}[1]{\textcolor[rgb]{0.00,0.23,0.31}{\textbf{#1}}}
\newcommand{\DataTypeTok}[1]{\textcolor[rgb]{0.68,0.00,0.00}{#1}}
\newcommand{\DecValTok}[1]{\textcolor[rgb]{0.68,0.00,0.00}{#1}}
\newcommand{\DocumentationTok}[1]{\textcolor[rgb]{0.37,0.37,0.37}{\textit{#1}}}
\newcommand{\ErrorTok}[1]{\textcolor[rgb]{0.68,0.00,0.00}{#1}}
\newcommand{\ExtensionTok}[1]{\textcolor[rgb]{0.00,0.23,0.31}{#1}}
\newcommand{\FloatTok}[1]{\textcolor[rgb]{0.68,0.00,0.00}{#1}}
\newcommand{\FunctionTok}[1]{\textcolor[rgb]{0.28,0.35,0.67}{#1}}
\newcommand{\ImportTok}[1]{\textcolor[rgb]{0.00,0.46,0.62}{#1}}
\newcommand{\InformationTok}[1]{\textcolor[rgb]{0.37,0.37,0.37}{#1}}
\newcommand{\KeywordTok}[1]{\textcolor[rgb]{0.00,0.23,0.31}{\textbf{#1}}}
\newcommand{\NormalTok}[1]{\textcolor[rgb]{0.00,0.23,0.31}{#1}}
\newcommand{\OperatorTok}[1]{\textcolor[rgb]{0.37,0.37,0.37}{#1}}
\newcommand{\OtherTok}[1]{\textcolor[rgb]{0.00,0.23,0.31}{#1}}
\newcommand{\PreprocessorTok}[1]{\textcolor[rgb]{0.68,0.00,0.00}{#1}}
\newcommand{\RegionMarkerTok}[1]{\textcolor[rgb]{0.00,0.23,0.31}{#1}}
\newcommand{\SpecialCharTok}[1]{\textcolor[rgb]{0.37,0.37,0.37}{#1}}
\newcommand{\SpecialStringTok}[1]{\textcolor[rgb]{0.13,0.47,0.30}{#1}}
\newcommand{\StringTok}[1]{\textcolor[rgb]{0.13,0.47,0.30}{#1}}
\newcommand{\VariableTok}[1]{\textcolor[rgb]{0.07,0.07,0.07}{#1}}
\newcommand{\VerbatimStringTok}[1]{\textcolor[rgb]{0.13,0.47,0.30}{#1}}
\newcommand{\WarningTok}[1]{\textcolor[rgb]{0.37,0.37,0.37}{\textit{#1}}}

\usepackage{longtable,booktabs,array}
\usepackage{calc} % for calculating minipage widths
% Correct order of tables after \paragraph or \subparagraph
\usepackage{etoolbox}
\makeatletter
\patchcmd\longtable{\par}{\if@noskipsec\mbox{}\fi\par}{}{}
\makeatother
% Allow footnotes in longtable head/foot
\IfFileExists{footnotehyper.sty}{\usepackage{footnotehyper}}{\usepackage{footnote}}
\makesavenoteenv{longtable}
\usepackage{graphicx}
\makeatletter
\newsavebox\pandoc@box
\newcommand*\pandocbounded[1]{% scales image to fit in text height/width
  \sbox\pandoc@box{#1}%
  \Gscale@div\@tempa{\textheight}{\dimexpr\ht\pandoc@box+\dp\pandoc@box\relax}%
  \Gscale@div\@tempb{\linewidth}{\wd\pandoc@box}%
  \ifdim\@tempb\p@<\@tempa\p@\let\@tempa\@tempb\fi% select the smaller of both
  \ifdim\@tempa\p@<\p@\scalebox{\@tempa}{\usebox\pandoc@box}%
  \else\usebox{\pandoc@box}%
  \fi%
}
% Set default figure placement to htbp
\def\fps@figure{htbp}
\makeatother
\usepackage{svg}





\setlength{\emergencystretch}{3em} % prevent overfull lines

\providecommand{\tightlist}{%
  \setlength{\itemsep}{0pt}\setlength{\parskip}{0pt}}



 


\KOMAoption{captions}{tableheading}
\makeatletter
\@ifpackageloaded{tcolorbox}{}{\usepackage[skins,breakable]{tcolorbox}}
\@ifpackageloaded{fontawesome5}{}{\usepackage{fontawesome5}}
\definecolor{quarto-callout-color}{HTML}{909090}
\definecolor{quarto-callout-note-color}{HTML}{0758E5}
\definecolor{quarto-callout-important-color}{HTML}{CC1914}
\definecolor{quarto-callout-warning-color}{HTML}{EB9113}
\definecolor{quarto-callout-tip-color}{HTML}{00A047}
\definecolor{quarto-callout-caution-color}{HTML}{FC5300}
\definecolor{quarto-callout-color-frame}{HTML}{acacac}
\definecolor{quarto-callout-note-color-frame}{HTML}{4582ec}
\definecolor{quarto-callout-important-color-frame}{HTML}{d9534f}
\definecolor{quarto-callout-warning-color-frame}{HTML}{f0ad4e}
\definecolor{quarto-callout-tip-color-frame}{HTML}{02b875}
\definecolor{quarto-callout-caution-color-frame}{HTML}{fd7e14}
\makeatother
\makeatletter
\@ifpackageloaded{bookmark}{}{\usepackage{bookmark}}
\makeatother
\makeatletter
\@ifpackageloaded{caption}{}{\usepackage{caption}}
\AtBeginDocument{%
\ifdefined\contentsname
  \renewcommand*\contentsname{Table of contents}
\else
  \newcommand\contentsname{Table of contents}
\fi
\ifdefined\listfigurename
  \renewcommand*\listfigurename{List of Figures}
\else
  \newcommand\listfigurename{List of Figures}
\fi
\ifdefined\listtablename
  \renewcommand*\listtablename{List of Tables}
\else
  \newcommand\listtablename{List of Tables}
\fi
\ifdefined\figurename
  \renewcommand*\figurename{Figure}
\else
  \newcommand\figurename{Figure}
\fi
\ifdefined\tablename
  \renewcommand*\tablename{Table}
\else
  \newcommand\tablename{Table}
\fi
}
\@ifpackageloaded{float}{}{\usepackage{float}}
\floatstyle{ruled}
\@ifundefined{c@chapter}{\newfloat{codelisting}{h}{lop}}{\newfloat{codelisting}{h}{lop}[chapter]}
\floatname{codelisting}{Listing}
\newcommand*\listoflistings{\listof{codelisting}{List of Listings}}
\makeatother
\makeatletter
\makeatother
\makeatletter
\@ifpackageloaded{caption}{}{\usepackage{caption}}
\@ifpackageloaded{subcaption}{}{\usepackage{subcaption}}
\makeatother
\usepackage{bookmark}
\IfFileExists{xurl.sty}{\usepackage{xurl}}{} % add URL line breaks if available
\urlstyle{same}
\hypersetup{
  pdftitle={Introduction to Computational Psychology},
  pdfauthor={Alan Jern},
  colorlinks=true,
  linkcolor={blue},
  filecolor={Maroon},
  citecolor={Blue},
  urlcolor={Blue},
  pdfcreator={LaTeX via pandoc}}


\title{Introduction to Computational Psychology}
\author{Alan Jern}
\date{2026-03-03}
\begin{document}
\maketitle

\renewcommand*\contentsname{Table of contents}
{
\hypersetup{linkcolor=}
\setcounter{tocdepth}{2}
\tableofcontents
}

\bookmarksetup{startatroot}

\chapter*{Welcome}\label{welcome}
\addcontentsline{toc}{chapter}{Welcome}

\markboth{Welcome}{Welcome}

\section*{Thanks}\label{thanks}
\addcontentsline{toc}{section}{Thanks}

\markright{Thanks}

This book borrows inspiration and content (especially Chapters
\hyperref[bayes]{2} and \hyperref[pragmatics]{7}) from Fausto Carcassi's
\href{https://github.com/thelogicalgrammar/cognitive_modelling_book}{Introduction
to Cognitive Modelling in R} book and I am tremendously grateful to him
for sharing his book publicly. Hopefully this resource will be equally
helpful to others.

A number of the modeling examples and homework assignments are adapted
from code written by Danielle Navarro from previous iterations of her
\href{https://compcogsci-3016.djnavarro.net/}{Computational Cognitive
Science course} with Andy Perfors. I am very grateful to both of them
for making all of their materials public (and well documented).

Your \href{mailto:jern@rose-hulman.edu}{feedback} is welcome and
encouraged.

\section*{License}\label{license}
\addcontentsline{toc}{section}{License}

\markright{License}

\href{https://creativecommons.org/licenses/by-nc/4.0/}{\pandocbounded{\includegraphics[keepaspectratio]{images/badges/cc-by-nc.png}}}

Anyone is free to reuse and adapt this book for their own non-commercial
purposes, with attribution. If you do use this book in any way,
\href{mailto:jern@rose-hulman.edu}{please tell me about it}.

\bookmarksetup{startatroot}

\chapter{Why computational modeling?}\label{intro}

The field of psychology is full of theories. But this book is focused on
a specific type of theory: the type that can be formalized as a
computational model.

What advantage does a computational (or mathematical) theory have over
one that can't be expressed in computational terms? The biggest
advantage is precision. For example, suppose you say that attention is
like a spotlight 🔦: we can attend to things that are currently within
the light, and we can control where the light is shining, but things
outside of the light are outside of our awareness. This is a kind of
theory -- an analogy-based one -- and it's a good start for making some
general \emph{qualitative} predictions about how attention works.

But you'll quickly run into problems if you want to make precise
\emph{quantitative} predictions about attention. How big is the
spotlight? Does the size expand and contract? How quickly does it move
around? Is attention completely absent outside of the spotlight or does
it ramp down as you get near the edge of the light? In other words, if
you wanted to build a computer model of this theory, a simple analogy
doesn't cut it.

\emph{Computational models force us to be explicit about our
assumptions.}

\section{Representations 🔸}\label{representations}

We don't perceive the world as it truly is. As one example, the visible
spectrum that our eyes can detect is just a fraction of the full
electromagnetic spectrum. Similarly, we can only hear a narrow range of
sound frequencies. In other words, we perceive an incomplete picture of
the world.

\begin{figure}[H]

{\centering \pandocbounded{\includegraphics[keepaspectratio]{images/01/range_of_sensation.png}}

}

\caption{Source: Abstruse Goose.}

\end{figure}%

Even for the parts of the world that we \emph{can} perceive, the
information we get through our senses is often incomplete, so we are
constantly making assumptions to fill in gaps.

As a result, what we have in our heads is a kind of model of the world
around us -- what cognitive scientists call a mental representation.
These representations help us to reach rapid conclusions about things
involving language, causes and effects, concepts, mental states, and
many other aspects of cognition.

Some of the key questions for cognitive scientists who use computational
models are:

\begin{enumerate}
\def\labelenumi{\arabic{enumi}.}
\tightlist
\item
  What mental representations do we rely on?
\item
  How do our minds use these representations to learn when we get new
  information?
\item
  What kind of information do we get and how do our expectations about
  the kind of information we're getting to affect how we use it?
\end{enumerate}

This book focuses on each of these questions.

\section{Homework 1: Build your first computational model
💻}\label{homework-1-build-your-first-computational-model}

To get some initial experience with computational modeling, you'll build
and experiment with a simple model of classical conditioning developed
by Robert Rescorla and Allan Wagner -- now called the Rescorla-Wagner
model.

All homework assignments for this book will be done in Google Colab.
Click the button below to view Homework 1.

\href{https://colab.research.google.com/drive/1tOD-ChXJ2MlsIqU85yWzSqZnmg64KcRv?usp=sharing}{\pandocbounded{\includesvg[keepaspectratio]{index_files/mediabag/colab-badge.svg}}}

\begin{tcolorbox}[enhanced jigsaw, opacitybacktitle=0.6, colback=white, title=\textcolor{quarto-callout-note-color}{\faInfo}\hspace{0.5em}{Note}, breakable, arc=.35mm, toptitle=1mm, rightrule=.15mm, toprule=.15mm, bottomtitle=1mm, coltitle=black, left=2mm, opacityback=0, colframe=quarto-callout-note-color-frame, leftrule=.75mm, bottomrule=.15mm, colbacktitle=quarto-callout-note-color!10!white, titlerule=0mm]

You'll have to make a copy of the notebook saved to your own Drive in
order to edit it.

\end{tcolorbox}

If you're unfamiliar with Colab (or Jupyter Notebooks), watch this brief
introduction video.

\url{https://www.youtube.com/embed/inN8seMm7UI}

This video gives some extra tips for working with Colab.

\url{https://youtu.be/rNgswRZ2C1Y?si=aVooywIqX0_G0NId}

\bookmarksetup{startatroot}

\chapter{Bayesian inference}\label{bayes}

This chapter reviews some basic probability and Bayesian inference. You
might be asking yourself: \emph{What does this have to do with
psychology?} The answer becomes clear when you recognize that most of
what we do when we're making sense of the world is drawing inferences.
When you see an ambiguous image like the one below, is it a rabbit or a
duck? When someone mumbles something, did they say ``hello'' or ``go to
hell''? When you take a pill and your headache goes away, did the pill
eliminate your headache or did the headache go away on its own?

\begin{figure}[H]

{\centering \includegraphics[width=0.7\linewidth,height=\textheight,keepaspectratio]{images/02/duck-rabbit.png}

}

\caption{The duck-rabbit illusion}

\end{figure}%

In all of these examples, there is more than one hypothesis about what
we observed. Probability and Bayesian inference provide the tools for
optimally determining how probable these different hypotheses are. One
of the claims of this book is that when people are making inferences in
situations like these, their inferences are often well predicted by the
optimal inferences dictated by probability theory.

\section{Basic probability 🎲}\label{basic-probability}

\begin{itemize}
\tightlist
\item
  Conditional probability: \(P(b|a) = \frac{P(a,b)}{P(a)}\)
\item
  Chain rule: \(P(a,b) = P(b|a)P(a)\)
\item
  Marginalization: \(P(d) = \sum_h P(d,h) = \sum_h P(d|h) P(h)\)
\item
  Bayes's rule:
  \(P(h|d) = \frac{P(d|h) P(h)}{P(d)} = \frac{P(d|h) P(h)}{\sum_h P(d|h) P(h)}\).
  \(P(d|h)\) is referred to as the \emph{likelihood}, \(P(h)\) is the
  \emph{prior}, and \(P(h|d)\) is the \emph{posterior}.
\end{itemize}

\section{A motivating example: Sampling from a bag
👝}\label{a-motivating-example-sampling-from-a-bag}

Suppose you have a bag full of black and red balls. You can't see inside
the bag and you don't know how many black and red balls are inside, but
you know that there are nine total balls in the bag.

You want to know how many black balls and red balls there are. There are
a finite number of hypotheses: \{0 black balls, 1 black ball, 2 black
balls, \ldots, 9 black balls\}. Let's call these hypotheses \(B_0\),
\(B_1\), etc., respectively.

You don't know which hypothesis is true, but you might have some idea
that some hypotheses are more likely than others. It's natural to
represent your uncertainty with a probability distribution over the
possible unknown states that the world could be in -- in this case, the
ten hypotheses. Each hypothesis gets assigned a probability and the
probabilities sum to 1.

For starters, let's assume that you \emph{don't} have any idea which
hypotheses are more likely. In other words, you give every hypothesis
the same probability: 1/10 = 0.1. This is called a uniform distribution
over hypotheses. This distribution is your prior.

Now suppose you put your hand in the bag and pull out a ball at random.
The possible observations are: \texttt{\{black,\ red\}}, Let's call them
\(B\) and \(R\), respectively. The probability of observing each color
depends on which hypothesis is true, i.e., how many balls of each color
are in the bag. For instance, if \(B_0\) is true (there are 0 black
balls in the bag), then the probability of observing a red ball is 1
(\(P(R|B_0)=1\)), and the probability of observing a black ball is 0
(\(P(B|B_0)=0\)). These expressions that tell us how probable our
observations are, given a specific hypothesis, are your likelihoods.

\subsection{Sampling from the generative
model}\label{sampling-from-the-generative-model}

Now we have a distribution over hypotheses (a prior), \(P(h)\), and a
distribution over observations given each hypothesis (a likelihood),
\(P(d|h)\). These two things allow us to create a \emph{generative
model}, a model for sampling new data.

How do we sample from the generative model? Note that which hypothesis
\(h\) is true does not depend on the data, while the data \(d\) depends
on which hypothesis is true. Therefore, we can sample from the
generative model using the following two-step process:

\begin{enumerate}
\def\labelenumi{\arabic{enumi}.}
\tightlist
\item
  Sample a hypothesis from the prior.
\item
  Sample data given the hypothesis, using the likelihood.
\end{enumerate}

Let's first create an array with the probability of each hypothesis:

\begin{Shaded}
\begin{Highlighting}[]
\ImportTok{import}\NormalTok{ numpy }\ImportTok{as}\NormalTok{ np}
\ImportTok{import}\NormalTok{ random}

\NormalTok{random.seed(}\DecValTok{2022}\NormalTok{) }\CommentTok{\# set random seed to get same results every time}

\NormalTok{h\_priors }\OperatorTok{=}\NormalTok{ np.repeat(}\FloatTok{0.1}\NormalTok{,}\DecValTok{10}\NormalTok{)}
\BuiltInTok{print}\NormalTok{(h\_priors)}
\end{Highlighting}
\end{Shaded}

Now we'll do the first step: create an array of 10000 hypotheses sampled
from the prior:

\begin{Shaded}
\begin{Highlighting}[]
\NormalTok{prior\_samples }\OperatorTok{=}\NormalTok{ np.array(random.choices(np.arange(}\DecValTok{0}\NormalTok{,}\DecValTok{10}\NormalTok{,}\DecValTok{1}\NormalTok{),}
\NormalTok{                   weights }\OperatorTok{=}\NormalTok{ h\_priors, }
\NormalTok{                   k }\OperatorTok{=} \DecValTok{10000}\NormalTok{))}
                   
\BuiltInTok{print}\NormalTok{(prior\_samples[}\DecValTok{0}\NormalTok{:}\DecValTok{9}\NormalTok{]) }\CommentTok{\# printing out just a few}
\end{Highlighting}
\end{Shaded}

Here, each number corresponds to one hypothesis: 0 corresponds to
\(B_0\), 1 to \(B_1\), and so on. Each sample represents one possible
way (a hypothesis) the world could be. Since the prior was uniform (each
hypothesis had the same probability), each hypothesis appears about
equally often. We can plot all the samples to verify:

\begin{Shaded}
\begin{Highlighting}[]
\ImportTok{import}\NormalTok{ seaborn }\ImportTok{as}\NormalTok{ sns}
\ImportTok{import}\NormalTok{ matplotlib.pyplot }\ImportTok{as}\NormalTok{ plt}

\NormalTok{sns.set\_theme()}
\NormalTok{fig, ax }\OperatorTok{=}\NormalTok{ plt.subplots()}
\NormalTok{sns.histplot(prior\_samples, bins}\OperatorTok{=}\DecValTok{10}\NormalTok{)}
\NormalTok{ax.set\_xlabel(}\StringTok{\textquotesingle{}Prior sample\textquotesingle{}}\NormalTok{)}
\NormalTok{ax.set\_ylabel(}\StringTok{\textquotesingle{}Number of samples\textquotesingle{}}\NormalTok{)}
\end{Highlighting}
\end{Shaded}

Now for the next step. For each sample in \texttt{prior\_samples}, we
want to sample an observation. To do that, let's pause for a second and
think about the probability of pulling a black ball given that
hypothesis \(B_3\) is true, for example. This means that there are 3
black balls and 6 red balls in the bag. So the probability of pulling a
black ball from the bag at random will be 3/9.

Generalizing this idea, we can get the probability of pulling a black
ball from the bag by dividing the elements of \texttt{prior\_samples} by
9:

\begin{Shaded}
\begin{Highlighting}[]
\NormalTok{p\_black }\OperatorTok{=}\NormalTok{ prior\_samples }\OperatorTok{/} \DecValTok{9}
\BuiltInTok{print}\NormalTok{(p\_black[}\DecValTok{0}\NormalTok{:}\DecValTok{4}\NormalTok{]) }\CommentTok{\# print out just a few}
\end{Highlighting}
\end{Shaded}

Now, to complete our generative model, we just need to sample one value
for each element of \texttt{p\_black}. Each sample represents a draw
from a bag.

\begin{Shaded}
\begin{Highlighting}[]
\NormalTok{ball\_samples }\OperatorTok{=}\NormalTok{ np.random.binomial(n }\OperatorTok{=}\NormalTok{ np.repeat(}\DecValTok{1}\NormalTok{,}\BuiltInTok{len}\NormalTok{(p\_black)),}
\NormalTok{                                  p }\OperatorTok{=}\NormalTok{ p\_black)}
\BuiltInTok{print}\NormalTok{(ball\_samples[}\DecValTok{0}\NormalTok{:}\DecValTok{9}\NormalTok{])}
\end{Highlighting}
\end{Shaded}

\texttt{ball\_samples} is 1 for black and 0 for red. Once again, let's
plot all our samples.

\begin{Shaded}
\begin{Highlighting}[]
\NormalTok{fig, ax }\OperatorTok{=}\NormalTok{ plt.subplots()}
\NormalTok{sns.histplot(ball\_samples, bins}\OperatorTok{=}\DecValTok{2}\NormalTok{)}
\NormalTok{ax.set\_xticks(ticks}\OperatorTok{=}\NormalTok{[}\DecValTok{0}\NormalTok{,}\DecValTok{1}\NormalTok{])}
\NormalTok{ax.set\_ylabel(}\StringTok{"Number of samples"}\NormalTok{)}
\end{Highlighting}
\end{Shaded}

You can think of this plot representing your overall beliefs about the
number of red and black balls in the bag, averaged over all possible
hypotheses.

Not surprisingly, we got about equal numbers of red and black balls.
This makes sense: We didn't have any prior expectations about whether
red or black balls were more likely in the bag.

How should your beliefs change after you pull a ball out of the bag?
That is, how should you respond to evidence?

\section{Bayesian updating: Learning from evidence
🤔}\label{bayesian-updating-learning-from-evidence}

Let's apply \hyperref[basic-probability]{Bayes's rule} to see how to
optimally incorporate new data into your beliefs.

\subsection{Applying Bayes's rule to the bag
case}\label{applying-bayess-rule-to-the-bag-case}

Suppose you have a uniform prior distribution over the 10 hypotheses
about balls in the bag. Now you pick a ball and it's black. Given this
observation \(B\), how should you change the probabilities you give to
each hypothesis?

Intuitively, you should now give a little bit more probability to those
hypotheses that have more black balls than red balls, because those are
the hypotheses that make your observations more likely. Moreover, you
can definitively exclude hypothesis \(B_0\), because your observation
would be impossible if \(B_0\) were true. Let's calculate this with
Bayes's rule.

The prior is the array \texttt{h\_priors} defined above. Given that we
have observed \(B\), the likelihood should tell us, for each hypothesis,
the probability of \(B\) given that hypothesis. For example, for
\(B_9\), the likelihood \(P(B|B_9) = 1\). For \(B_8\),
\(P(B|B_8) = 8/9\), because 8 of the 9 balls are black.

Generalizing this idea, \(P(B|B_n) = n/9\). We can therefore compute the
likelihoods for all hypotheses in a vector:

\begin{Shaded}
\begin{Highlighting}[]
\NormalTok{likelihoods }\OperatorTok{=}\NormalTok{ np.arange(}\DecValTok{0}\NormalTok{,}\DecValTok{10}\NormalTok{,}\DecValTok{1}\NormalTok{) }\OperatorTok{/} \DecValTok{9}

\BuiltInTok{print}\NormalTok{(likelihoods)}
\end{Highlighting}
\end{Shaded}

Now suppose we want to find the probability of hypothesis \(B_5\) after
observing one draw \(B\). Let's apply Bayes's rule:

\[P(B_5 | B) = \frac{P(B|B_5) P(B_5)}{\sum_h{p(B|h) P(h)}}\]

Let's compute the parts we need to calculate \(P(B_5 | B)\).

\begin{Shaded}
\begin{Highlighting}[]
\CommentTok{\# Prior}
\NormalTok{p\_B5 }\OperatorTok{=}\NormalTok{ h\_priors[}\DecValTok{3}\NormalTok{]}

\CommentTok{\# Likelihood}
\NormalTok{likelihood\_B5 }\OperatorTok{=}\NormalTok{ likelihoods[}\DecValTok{5}\NormalTok{]}

\CommentTok{\# Data}
\NormalTok{p\_B }\OperatorTok{=} \BuiltInTok{sum}\NormalTok{(likelihoods}\OperatorTok{*}\NormalTok{h\_priors)}

\CommentTok{\# Posterior}
\NormalTok{p\_B5\_given\_B }\OperatorTok{=}\NormalTok{ p\_B5 }\OperatorTok{*}\NormalTok{ likelihood\_B5 }\OperatorTok{/}\NormalTok{ p\_B}

\CommentTok{\# Print out results}
\BuiltInTok{print}\NormalTok{(}\StringTok{"P(B5) = "} \OperatorTok{+} \BuiltInTok{str}\NormalTok{(p\_B5)) }\CommentTok{\# Prior}
\BuiltInTok{print}\NormalTok{(}\StringTok{"P(B|B5) = "} \OperatorTok{+} \BuiltInTok{str}\NormalTok{(likelihood\_B5)) }\CommentTok{\# Likelihood}
\BuiltInTok{print}\NormalTok{(}\StringTok{"P(B) = "} \OperatorTok{+} \BuiltInTok{str}\NormalTok{(p\_B)) }\CommentTok{\# Data}
\BuiltInTok{print}\NormalTok{(}\StringTok{"P(B5|B) = "} \OperatorTok{+} \BuiltInTok{str}\NormalTok{(p\_B5\_given\_B)) }\CommentTok{\# Posterior}
\end{Highlighting}
\end{Shaded}

Let's update the probabilities for all hypotheses in a more compact way.

\begin{Shaded}
\begin{Highlighting}[]
\NormalTok{posteriors }\OperatorTok{=}\NormalTok{ (likelihoods }\OperatorTok{*}\NormalTok{ h\_priors) }\OperatorTok{/} \BuiltInTok{sum}\NormalTok{(likelihoods }\OperatorTok{*}\NormalTok{ h\_priors)}

\ControlFlowTok{for}\NormalTok{ i }\KeywordTok{in} \BuiltInTok{range}\NormalTok{(}\BuiltInTok{len}\NormalTok{(posteriors)):}
  \BuiltInTok{print}\NormalTok{(}\StringTok{"P(B"} \OperatorTok{+} \BuiltInTok{str}\NormalTok{(i) }\OperatorTok{+} \StringTok{"|B) = "} \OperatorTok{+} \BuiltInTok{str}\NormalTok{(posteriors[i]))}
\end{Highlighting}
\end{Shaded}

As expected, Bayes's rule says we should increase the probability we
assign to hypotheses with more black balls than red balls. Additionally,
let's double-check that the posterior probabilities sum to 1 (a
requirement for a valid probability distribution).

\begin{Shaded}
\begin{Highlighting}[]
\BuiltInTok{sum}\NormalTok{(posteriors)}
\end{Highlighting}
\end{Shaded}

Finally, let's plot the posterior probabilities.

\begin{Shaded}
\begin{Highlighting}[]
\NormalTok{hypotheses }\OperatorTok{=}\NormalTok{ (}\StringTok{\textquotesingle{}B0\textquotesingle{}}\NormalTok{, }\StringTok{\textquotesingle{}B1\textquotesingle{}}\NormalTok{, }\StringTok{\textquotesingle{}B2\textquotesingle{}}\NormalTok{, }\StringTok{\textquotesingle{}B3\textquotesingle{}}\NormalTok{, }\StringTok{\textquotesingle{}B4\textquotesingle{}}\NormalTok{,}
              \StringTok{\textquotesingle{}B5\textquotesingle{}}\NormalTok{, }\StringTok{\textquotesingle{}B6\textquotesingle{}}\NormalTok{, }\StringTok{\textquotesingle{}B7\textquotesingle{}}\NormalTok{, }\StringTok{\textquotesingle{}B8\textquotesingle{}}\NormalTok{, }\StringTok{\textquotesingle{}B9\textquotesingle{}}\NormalTok{)}
\NormalTok{y\_pos }\OperatorTok{=}\NormalTok{ np.arange(}\BuiltInTok{len}\NormalTok{(hypotheses))}

\NormalTok{fig, ax }\OperatorTok{=}\NormalTok{ plt.subplots()}
\NormalTok{sns.barplot(posteriors, orient}\OperatorTok{=}\StringTok{"y"}\NormalTok{)}
\NormalTok{ax.set\_yticks(ticks }\OperatorTok{=}\NormalTok{ y\_pos, labels}\OperatorTok{=}\NormalTok{hypotheses)}
\NormalTok{ax.set\_ylabel(}\StringTok{\textquotesingle{}Hypothesis\textquotesingle{}}\NormalTok{)}
\NormalTok{ax.set\_xlabel(}\StringTok{\textquotesingle{}Probability\textquotesingle{}}\NormalTok{)}
\end{Highlighting}
\end{Shaded}

\subsection{How to avoid calculating P(d)}\label{normalization}

In practice, we generally do not need to calculate the \(P(d)\) (the
denominator in Bayes's rule) explicitly. I'll give you the general idea
why in this section.

First, we create an array of prior probabilities, which has as many
components as there are hypotheses. We'll just reuse \texttt{h\_priors}.
Note that the probabilities sum to 1, as they should because it's a
probability distribution.

\begin{Shaded}
\begin{Highlighting}[]
\NormalTok{h\_priors}
\end{Highlighting}
\end{Shaded}

Next, we create a likelihood array. When we did calculations above, we
only had an array with the likelihoods for a specific observation.
However, we would like to have something that encodes the likelihood
function for each possible observation given each possible hypothesis,
rather than just for a specific observation.

In this example, there are two possible observations: \(B\) and \(R\).
We can encode the likelihood as an \(m \times n\) array where \(m\) is
the number of hypotheses and \(n\) is the number of possible
observations. In our case: \(10 \times 2\).

\begin{Shaded}
\begin{Highlighting}[]
\NormalTok{likelihood\_array }\OperatorTok{=}\NormalTok{ np.array((np.arange(}\DecValTok{0}\NormalTok{,}\DecValTok{10}\NormalTok{,}\DecValTok{1}\NormalTok{) }\OperatorTok{/} \DecValTok{9}\NormalTok{,}
                             \DecValTok{1}\OperatorTok{{-}}\NormalTok{(np.arange(}\DecValTok{0}\NormalTok{,}\DecValTok{10}\NormalTok{,}\DecValTok{1}\NormalTok{) }\OperatorTok{/} \DecValTok{9}\NormalTok{))).T}
\BuiltInTok{print}\NormalTok{(likelihood\_array)}
\end{Highlighting}
\end{Shaded}

Now we multiply the prior and likelihoods together (the numerator of
Bayes's rule) \emph{element-wise} (first element gets multiplied with
first element, second element by second element, etc.):

\begin{Shaded}
\begin{Highlighting}[]
\NormalTok{prior\_array }\OperatorTok{=}\NormalTok{ np.array((h\_priors, h\_priors)).T}
\NormalTok{bayes\_numerator }\OperatorTok{=}\NormalTok{ likelihood\_array }\OperatorTok{*}\NormalTok{ prior\_array}

\BuiltInTok{print}\NormalTok{(bayes\_numerator)}
\end{Highlighting}
\end{Shaded}

Finally, we want a distribution for each column, i.e., a distribution
over hypotheses given each observation. Therefore, we sum each column
and then divide each element by the sum of its column:

\begin{Shaded}
\begin{Highlighting}[]
\NormalTok{posteriors }\OperatorTok{=}\NormalTok{ bayes\_numerator }\OperatorTok{/}\NormalTok{ np.}\BuiltInTok{sum}\NormalTok{(bayes\_numerator, axis }\OperatorTok{=} \DecValTok{0}\NormalTok{)}
\BuiltInTok{print}\NormalTok{(posteriors)}
\end{Highlighting}
\end{Shaded}

And that gives us the posterior without us having to explicitly
calculate the evidence for each observation!

\begin{tcolorbox}[enhanced jigsaw, opacitybacktitle=0.6, colback=white, title=\textcolor{quarto-callout-note-color}{\faInfo}\hspace{0.5em}{Note}, breakable, arc=.35mm, toptitle=1mm, rightrule=.15mm, toprule=.15mm, bottomtitle=1mm, coltitle=black, left=2mm, opacityback=0, colframe=quarto-callout-note-color-frame, leftrule=.75mm, bottomrule=.15mm, colbacktitle=quarto-callout-note-color!10!white, titlerule=0mm]

The general idea is this. Because the denominator of Bayes's rule, for a
fixed observation, is a constant, you can usually get away with
computing \(P(d|h) P(h)\) for every possible hypothesis \(h\) and then
``normalize'' the resulting values so that they sum to 1 (remember that
they have to in order for it to be a valid probability distribution).

\end{tcolorbox}

\section{Exercises 📝}\label{bayes-exercises}

\subsection{Taxi cabs}\label{taxi-cabs}

80\% of the taxi cabs in Simpletown are green and 20\% are
yellow\footnote{After the protracted and aptly named ``color wars'' of
  2113-2114, the previously dominant (and ruthless) orange cabs were
  wiped from existence. For the purposes of this problem, you may ignore
  them.}. A hit-and-run accident happened at night involving a taxi. A
witness claimed that the taxi was yellow. After extensive testing, it is
determined that the witness can correctly identify the color of a taxi
only 75\% of the time under conditions like the ones present during the
accident. What is the probability that the taxi was yellow?

\subsection{Flipping coins}\label{flipping-coins}

You observe a sequence of coin flips and want to determine if the coin
is a trick coin (always comes up heads) or a normal coin. Let
\(P(\text{heads}) = \theta\). Let \(h_1\) be the hypothesis that
\(\theta = 0.5\) (fair coin). Let \(h_2\) be the hypothesis that
\(\theta = 1\) (trick coin).

For this problem, we will define something called prior odds, which is
the ratio of prior probabilities assigned to two hypotheses:
\(\frac{P(h_1)}{P(h_2)}\). Because most coins aren't trick coins, we
assume that \(\frac{P(h_1)}{P(h_2)} = 999\), indicating a very strong
(999 to 1) prior probability in favor of fair coins. We can now compute
the posterior odds, the ratio of posterior probabilities for the two
hypotheses after observing some data \(d\):
\(\frac{P(h_1|d)}{P(h_2|d)}\).

Compute the posterior odds after observing the following sequences of
coin flips:

\begin{enumerate}
\def\labelenumi{\arabic{enumi}.}
\tightlist
\item
  HHTHT
\item
  HHHHH
\item
  HHHHHHHHHH
\end{enumerate}

\section{Solutions}\label{solutions}

\subsection{Taxi cabs}\label{taxi-cabs-1}

Let \(h_1\) be the hypothesis that the taxi is yellow. Let \(h_2\) be
the hypothesis that the taxi is green. Let data \(d\) be the witness
report that the taxi was yellow. Given the problem statement,
\(P(h_1) = 0.2\) and \(P(h_2) = 0.8\). The witness is only accurate 75\%
of the time, so \(P(d|h1) = 0.75\) (the witness saw a yellow taxi and
correctly identified it) and \(P(d|h2) = 0.25\) (the witness saw a green
taxi but identified it as yellow). Now we apply Bayes's rule:

\[\begin{align}
P(h_1|d) &= \frac{P(d|h_1) P(h_1)}{P(d)} \\
&= \frac{P(d|h_1) P(h_1)}{P(d|h_1) P(h_1) + P(d|h_2) P(h_2)} \\
&= \frac{(0.75) (0.2)}{(0.75)(0.2) + (0.25)(0.8)} \approx 0.43
\end{align}\]

Because yellow cabs are rare (have low prior probability), it is
actually more probable that the cab was green, even though the witness
is 75\% accurate.

\subsection{Flipping coins}\label{flipping-coins-1}

\subsubsection*{HHTHT}\label{hhtht}
\addcontentsline{toc}{subsubsection}{HHTHT}

\[
\begin{align}
\frac{P(h_1|d)}{P(h_2|d)} &= \frac{P(d|h_1)}{P(d|h_2)} \frac{P(h_1)}{P(h_2)} \\
&= \frac{(1/2)^5}{0} \times 999 = \inf
\end{align}
\] This sequence isn't even possible under \(h_2\) so we have infinite
evidence in favor of \(h_1\).

\subsubsection*{HHHHH}\label{hhhhh}
\addcontentsline{toc}{subsubsection}{HHHHH}

\[
\begin{align}
\frac{P(h_1|d)}{P(h_2|d)} &= \frac{P(d|h_1)}{P(d|h_2)} \frac{P(h_1)}{P(h_2)} \\
&= \frac{(1/2)^5}{1^5} \times 999 = 31.2
\end{align}
\]

This sequence favors \(h_1\) by a factor of about 31. Even five heads in
a row can't overcome our strong prior favoring \(h_1\).

\subsubsection*{HHHHHHHHHH}\label{hhhhhhhhhh}
\addcontentsline{toc}{subsubsection}{HHHHHHHHHH}

\[
\begin{align}
\frac{P(h_1|d)}{P(h_2|d)} &= \frac{P(d|h_1)}{P(d|h_2)} \frac{P(h_1)}{P(h_2)} \\
&= \frac{(1/2)^{10}}{1^{10}} \times 999 = 0.98
\end{align}
\]

Now the evidence favors \(h_2\) (trick coin) just barely.

\bookmarksetup{startatroot}

\chapter{Generalization}\label{generalization}

In the last chapter, you learned that much of cognition is about making
inferences. A common inference we're faced with involves generalizing
examples of things to new cases.

\begin{itemize}
\tightlist
\item
  A child hears a brand new word and has to figure out which objects
  that word applies to.
\item
  You eat one candy from an assorted box and then try to guess what the
  others might taste like.
\item
  You remember that your friend liked doing crosswords during a weekend
  trip at the cabin one time, so you guess they might like a book of
  puzzles as a gift (generalizing their interests).
\end{itemize}

How can we apply Bayesian inference to these kinds of problems?

\section{Healthy hormone levels 💉}\label{hormones}

The following example comes from a
\href{https://cocosci.princeton.edu/tom/papers/TGbbs.pdf}{2001 paper by
Josh Tenenbaum and Thomas Griffiths}\footnote{Tenenbaum, J. B. \&
  Griffiths, T. L. (2001). Generalization, similarity, and Bayesian
  inference. Behavioral and Brain Sciences, 24, 629-640.}

\textbf{The basic problem}: You learn the value of a healthy hormone
level (say, 60) that varies on a scale from 1 to 100 (integers only).
What is the probability that another value (say, 70) is also healthy?

\subsection{Setting up a model}\label{setting-up-a-model}

\subsubsection{The hypothesis space}\label{the-hypothesis-space}

To start with, we'll assume that healthy values lie in a contiguous
interval. Using the term from the paper, this interval is the
\emph{consequential region} \(C\).

The hypothesis space consists of all possible consequential regions. For
example, {[}0,100{]}, {[}10,19{]}, and {[}44,45{]}, are all valid
hypotheses. The full hypothesis space is every valid interval between 0
and 100.

\subsubsection{Prior}\label{prior}

How much weight should we assign to each hypothesis? We might have
reason to favor shorter intervals over longer ones, for example. In the
paper, they use an
\href{https://en.wikipedia.org/wiki/Erlang_distribution}{Erlang prior}.
Alternatively, for simplicity of calculation, we could again assume a
uniform prior distribution, placing equal weight on all hypotheses, like
we did in the previous chapter. This is tantamount to making no prior
assumptions about which intervals are most probable.

\subsubsection{Likelihood}\label{likelihood}

Suppose you learn that a healthy patient has a hormone level of 60. What
was the likelihood of observing this value, assuming we know which
hypothesis is correct? That is, what is \(P(x = 60 | h)\). It depends on
how we assume the patient was chosen.

\paragraph{Weak vs.~strong sampling}\label{weak-strong-sampling}

Under \emph{weak sampling}, we assume that each observation was sampled
from the full range of possibilities, and it was just a coincidence that
we happened to get one from the consequential region (a healthy
patient). If that's true, then the probability of getting any particular
value doesn't depend on which hypothesis is true:

\[
P(x|h) = \frac{1}{L}
\]

where \(L\) is the length of the range of possible values (100 in our
case).

Under \emph{strong sampling}, we assume that each observation was
specifically chosen as an example of the consequential region \(C\). In
other words, someone chose a healthy person and tested their hormone
levels as an example for you. In this case, the probability of seeing a
particular value depends on the size of the region:

\[
P(x|h) = \begin{cases} 
  \frac{1}{|h|} & \text{if } x \in h \\
  0 & \text{otherwise}
  \end{cases}
\] where \(|h|\) is the \emph{size} of \(h\), i.e., the number of values
contained in \(h\). If you have multiple observations
\(X = \{x_1, x_2, \ldots, x_n \}\), then \(P(X|h) = (1/|h|)^n\). This is
because we will assume each sample is independent like a coin flip.

The result of the strong sampling assumption is the \emph{size
principle}: among hypotheses that include all of the observed examples,
those that are smallest will receive higher posterior probability
because they will have higher likelihoods.

\begin{figure}[H]

{\centering \includegraphics[width=0.9\linewidth,height=\textheight,keepaspectratio]{images/03/hypotheses.png}

}

\caption{Sample hypotheses. The thickness of the lines indicates their
likelihood, depicting the size principle. Image from Griffiths \&
Tenenbaum (2001).}

\end{figure}%

\subsubsection{Posterior}\label{posterior}

We can now simply apply Bayes's rule to compute the probability of each
hypothesis, given an observation (or set of observations).

\[
P(h|X) = \frac{P(X|h) P(h)}{\sum_{h_i} P(X|h_i) P(h_i)}
\]

\subsection{Generalizing}\label{generalizing}

We aren't quite finished. Remember that we really want to know the
probability of some new value \(y\) also being a healthy hormone level.
But at this point all we have done is assign a probability to each
\emph{interval} being the consequential region.

What we want to do is essentially a two-step process:

\begin{enumerate}
\def\labelenumi{\arabic{enumi}.}
\tightlist
\item
  For each hypothesis \(h\), check to see if \(y\) is in it.
\item
  If it is, check how probable it is that \(h\) is the consequential
  region \(C\), given our observations \(X\).
\end{enumerate}

This basic idea is sometimes known as hypothesis averaging because we
don't actually care which hypothesis is the right one, so we'll just
average over \emph{all} hypotheses, weighted by how probable they are.
Specifically, we'll compute:

\[
P(y \in C|X) = \sum_h P(y \in C | h) P(h | X)
\]

The second term on the right is what we computed earlier using Bayes's
rule.

What about the first term? This time, we'll assume \emph{weak sampling}
because there's no reason to assume that this new value \(y\) was chosen
to be a healthy value or not.

\subsection{Homework 2: Finish the
details}\label{homework-2-finish-the-details}

\href{https://colab.research.google.com/drive/1aqPpMyc-eMyLGpHFl2Fq_NlLiTuHkGky?usp=sharing}{\pandocbounded{\includesvg[keepaspectratio]{index_files/mediabag/colab-badge.svg}}}

I haven't provided all the details for this model because your
assignment is to finish the implementation yourself, run some
simulations, and collect a small amount of real data to compare the
model to.

\section{The number game 🔢}\label{the-number-game}

In most domains, requiring concepts to be restricted to contiguous
intervals is not realistic. Numbers are one example. Consider the space
of possible number concepts you could make up for integers between 1 to
100. In addition to concepts like ``numbers between 20 and 50'', there
are many other plausible concepts like ``multiples of 10'', ``even
numbers'', or ``powers of 3''.

Consider the following problem: You are given one or more examples \(X\)
of numbers that fit some rule and you want to know how probable it is
that a new number \(y\) also fits the rule.

The model we discussed before can be naturally extended to this problem.
For the likelihood, we can make the same
\hyperref[weak-strong-sampling]{strong sampling} assumption as before.

The prior is where things get a little trickier. Intuitively, some
concepts like ``even numbers'' seem more probable than concepts like
``multiples of 7''. This now becomes a psychological question: Which
rules will people find to be more intuitively plausible? There is no
single way to decide this, but we could run a survey to find out: Give
people a long list of rules and ask them to judge how intuitively
natural they seem. We could then construct a prior probability
distribution using this data.

Alternatively, we could come up with some definition of ``complexity''
in hypotheses and assume that less complex hypotheses will receive
higher prior probability.

Once we have chosen a prior probability distribution \(P(h)\), we can
now proceed just as we did before.

\section{Inductive generalizations about animal properties
🐴}\label{inductive-generalizations-about-animal-properties}

Now let's consider an even more complex generalization problem, based on
a \href{http://sanjanalab.org/reprints/Sanjana_NIPS_2002.pdf}{2002 paper
by Neville Sanjana and Josh Tenenbaum}\footnote{Sanjana, N. \&
  Tenenbaum, J. (2002). Bayesian models of inductive generalization.
  Advances in Neural Information Processing Systems 15}. This is the
problem of generalizing properties from a set of example animals to
other animals. The paper uses the following example:

\begin{verbatim}
Chimps have blicketitis
Squirrels have blicketitis
--------------------------
Horses have blicketitis
\end{verbatim}

The way to read this is as follows: The premises state that
\textbf{chimps and squirrels have blicketitis}. The conclusion is that
\textbf{horses have blicketitis}. The inductive generalization question
is how probable is the conclusion given the premises? Intuitively, the
conclusion in this example seems more plausible than the conclusion in
the following example:

\begin{verbatim}
Chimps have blicketitis
Gorillas have blicketitis
--------------------------
Horses have blicketitis
\end{verbatim}

The interesting psychological question is why is it that some
generalizations seem more intuitively plausible than others?

\subsection{Hypothesis space}\label{hypothesis-space}

This problem is conceptually similar to the ones we've already been
discussing. You want to infer what animals have blicketitis after seeing
some examples of animals that blicketitis. The first question to answer
is: What is the analogue of a consequential region for this problem?
Animals don't naturally fall on a one-dimensional interval so we'll need
to define a different hypothesis space. One possibility is a hierarchy,
which naturally captures the knowledge people have about animals.

\subsubsection{Clustering}\label{clustering}

The paper first creates a hierarchy of eight animals using similarity
data collected from people. Specifically, they asked people to judge how
similar all pairs of eight animals were and then calculated the average
similarity judgment for each animal.

These similarity judgments can be used to construct a tree using a
simple clustering algorithm. The algorithm works as follows:

\begin{enumerate}
\def\labelenumi{\arabic{enumi}.}
\tightlist
\item
  Put all animals in their own cluster.
\item
  Where there is more than one cluster that hasn't been placed in a
  group, do the following:

  \begin{enumerate}
  \def\labelenumii{\arabic{enumii}.}
  \tightlist
  \item
    Identify the pair of clusters with the greatest similarity between
    them.
  \item
    Group those clusters into their own new cluster.
  \end{enumerate}
\end{enumerate}

There are several approaches for computing the similarity between two
clusters that contain multiple animals. For example, you might use the
maximum similarity between any pair of individual animals in the two
clusters.

The results of this algorithm can be represented as a tree, shown below.
Each node in the tree represents a cluster. The hypotheses we will
consider will be any combination of 1, 2, or 3 of the clusters
determined using the clustering algorithm.

\begin{figure}[H]

{\centering \includegraphics[width=0.9\linewidth,height=\textheight,keepaspectratio]{images/03/animal_clusters.png}

}

\caption{The tree of animal species. Image from Sanjana \& Tenenbaum
(2002).}

\end{figure}%

\subsection{The model}\label{the-model}

We can now define the model. First, let's define \(P(h)\) where \(h\) is
a set of clusters. The authors make an assumption analogous to the
following:

\[
P(h) \propto \frac{1}{\phi^k}
\] where \(k\) is the number of clusters in \(h\). \(\phi\) is a
parameter that we can choose. As long as \(\phi > 1\), \(P(h)\) will be
smaller for hypotheses consisting of more clusters. This has the effect
of assigning higher weight to ``simpler'' hypotheses.

Once again, we will make the \hyperref[sampling-assumptions]{strong
sampling assumption} for the likelihood. This time, \(|h|\) is the
number of animal species in \(h\) and \(n\) is the number of examples in
the premises.

A consequence of both of these assumptions is something like
\href{https://en.wikipedia.org/wiki/Occam\%27s_razor}{Occam's razor}
which says that the simplest explanation should be preferred. This model
will assign 0 likelihood to any hypotheses that don't include \(X\),
thus narrowing the hypothesis space down to just hypotheses that are
possible. Of those, it will favor hypotheses with fewer clusters and
with fewer animals. In other words, it will favor the simplest
hypotheses that are consistent with the examples we've seen.

\subsection{Try out the model
yourself}\label{try-out-the-model-yourself}

You can try out a running version of the model by making a copy of the
\href{https://colab.research.google.com/drive/1QIhYLYMcNXjrRQBXKRz9CZ4xMirzqcsp?usp=sharing}{code
here}. Here, let's look at how the model handles a few specific cases.

Let's start with a single example: horses can get blicketitis.

\begin{Shaded}
\begin{Highlighting}[]
\NormalTok{animalGeneralization([}\StringTok{"horse"}\NormalTok{])}
\end{Highlighting}
\end{Shaded}

\pandocbounded{\includegraphics[keepaspectratio]{03-generalization_files/figure-pdf/cell-3-output-1.pdf}}

\begin{verbatim}
array([1.        , 0.52984834, 0.30628906, 0.30628906, 0.19579428,
       0.19579428, 0.12893614, 0.12893614, 0.0842413 , 0.0842413 ])
\end{verbatim}

Here, we see a standard generalization curve. Now, let's add a few more
animals.

\begin{Shaded}
\begin{Highlighting}[]
\NormalTok{animalGeneralization([}\StringTok{"horse"}\NormalTok{, }\StringTok{"cow"}\NormalTok{, }\StringTok{"mouse"}\NormalTok{])}
\end{Highlighting}
\end{Shaded}

\pandocbounded{\includegraphics[keepaspectratio]{03-generalization_files/figure-pdf/cell-4-output-1.pdf}}

\begin{verbatim}
array([1.        , 1.        , 0.57379051, 0.57379051, 0.47852518,
       0.47852518, 1.        , 0.60980466, 0.1707609 , 0.1707609 ])
\end{verbatim}

Now the model has increased probability for most other animals. This
makes sense because there is more reason to think blicketitis might
affect lots of different animals.

\begin{Shaded}
\begin{Highlighting}[]
\NormalTok{animalGeneralization([}\StringTok{"horse"}\NormalTok{, }\StringTok{"cow"}\NormalTok{, }\StringTok{"mouse"}\NormalTok{, }\StringTok{"squirrel"}\NormalTok{])}
\end{Highlighting}
\end{Shaded}

\pandocbounded{\includegraphics[keepaspectratio]{03-generalization_files/figure-pdf/cell-5-output-1.pdf}}

\begin{verbatim}
array([1.        , 1.        , 0.64983602, 0.64983602, 0.58531333,
       0.58531333, 1.        , 1.        , 0.18405475, 0.18405475])
\end{verbatim}

Adding a squirrel further supported this idea. But what if we add
additional examples of animals we've already seen?

\begin{Shaded}
\begin{Highlighting}[]
\NormalTok{animalGeneralization([}\StringTok{"horse"}\NormalTok{, }\StringTok{"cow"}\NormalTok{, }\StringTok{"mouse"}\NormalTok{, }\StringTok{"squirrel"}\NormalTok{, }\StringTok{"horse"}\NormalTok{, }\StringTok{"squirrel"}\NormalTok{])}
\end{Highlighting}
\end{Shaded}

\pandocbounded{\includegraphics[keepaspectratio]{03-generalization_files/figure-pdf/cell-6-output-1.pdf}}

\begin{verbatim}
array([1.        , 1.        , 0.32551484, 0.32551484, 0.26938358,
       0.26938358, 1.        , 1.        , 0.06883892, 0.06883892])
\end{verbatim}

Now the probabilities for other animals drop, because it's starting to
look like maybe this disease only affects the four animals we've seen so
far.

To sum up, when we've seen a small number of examples, like a single
horse, the model will generally prefer simpler hypotheses. But once
we've seen more data, it will favor more complex hypotheses (like a
group of animals from two separate evolutionary clusters) if the data
support it.

As one final example, let's look at the impact of multiple examples of a
single animal.

\begin{Shaded}
\begin{Highlighting}[]
\NormalTok{animalGeneralization([}\StringTok{"gorilla"}\NormalTok{])}
\NormalTok{animalGeneralization([}\StringTok{"gorilla"}\NormalTok{, }\StringTok{"gorilla"}\NormalTok{])}
\NormalTok{animalGeneralization([}\StringTok{"gorilla"}\NormalTok{, }\StringTok{"gorilla"}\NormalTok{, }\StringTok{"gorilla"}\NormalTok{])}
\end{Highlighting}
\end{Shaded}

\pandocbounded{\includegraphics[keepaspectratio]{03-generalization_files/figure-pdf/cell-7-output-1.pdf}}

\pandocbounded{\includegraphics[keepaspectratio]{03-generalization_files/figure-pdf/cell-7-output-2.pdf}}

\pandocbounded{\includegraphics[keepaspectratio]{03-generalization_files/figure-pdf/cell-7-output-3.pdf}}

\begin{verbatim}
array([0.01730482, 0.01730482, 0.01730482, 0.01730482, 0.1259028 ,
       1.        , 0.01270155, 0.01270155, 0.01111944, 0.01111944])
\end{verbatim}

Here, we see that the model becomes increasingly confident that the
property is unique to gorillas. This makes intuitive sense and it's
something that people seem to exhibit in their judgments. But, as they
point out in the paper, it's not something that non-probabilistic models
can easily explain.

\subsection{Results}\label{results}

The results in the paper show that this model predicts people's
judgments quite well, better than alternative models that do not rely on
Bayesian inference. These results suggest that the assumptions of the
model are very similar to the assumptions that people make when making
inductive generalizations.

\begin{figure}[H]

{\centering \includegraphics[width=0.9\linewidth,height=\textheight,keepaspectratio]{images/03/st_results.png}

}

\caption{Comparison between model results and human judgments. Image
from Sanjana \& Tenenbaum (2002).}

\end{figure}%

\bookmarksetup{startatroot}

\chapter{Categorization}\label{categorization}

The last chapter introduced the problem of inductive generalization.
This chapter will focus on a specific case of generalization that is of
particular interest to psychologists, cognitive scientists, and people
who use AI and machine learning: \emph{classification} or
\emph{categorization}. Basically, assigning labels to things.

People constantly classify things in the world into categories (chairs,
cats, friends, enemies, edible stuff, and so on). Doing this helps us to
communicate and function in novel situations.

In machine learning, it's often useful to classify inputs into different
categories like whether a social media post violates community standards
or not, whether an image contains a human face or not, or whether an MRI
contains a tumor.

Psychologists study how people form categories because it provides a
window into how we organize our knowledge. But the basic problem is the
same as the one studied in machine learning.

I'm going to introduce a fairly simple psychological model of
categorization introduced by Robert Nosofsky called the Generalized
Context Model. It assumes that people make classification judgments
using the following algorithm:

\begin{enumerate}
\def\labelenumi{\arabic{enumi}.}
\tightlist
\item
  Remember all the examples of categories you've seen before.
\item
  When you want to classify a new instance, compare it to all previous
  examples of different categories and rate how similar it is to each
  one.
\item
  Assign it to the category that has the highest average similarity.
\end{enumerate}

This model makes some outlandish assumptions (like the idea that people
would remember every example they had seen before). But, to a first
approximation, it does a decent job of predicting people's
classification judgments in a lot of situations. And for learning
purposes, it has the advantage of being pretty easy to understand and
implement.

\section{A typical category learning experiment
🟢🟨}\label{a-typical-category-learning-experiment}

In psychology, category learning experiments have a pretty similar
structure. People see some unfamiliar stimuli that differ on several
dimensions and can be sorted into different categories. The task is to
learn what distinguishes one category from another. The experiments
usually consist of two phases: a training phase, and a testing phase.

\begin{itemize}
\tightlist
\item
  \emph{Training phase}: People see many examples of each category and
  have to classify them, often simply by guessing at first. They get
  feedback and gradually learn to tell the different categories apart.
\item
  \emph{Testing phase}: People see more examples, usually some brand new
  ones, and have to classify them again. This time they don't get
  feedback. The point of this phase is to test what people actually
  learned about the meaning of the categories.
\end{itemize}

The stimuli could be abstract shapes, cartoon insects, race cars with
different features, or really anything. What matters is that they have
clearly distinguishable features.

\section{Representing stimuli in a model
🧩}\label{representing-stimuli-in-a-model}

For modeling purposes, all of these types of stimuli (with
discrete-valued features) can be represented in a matrix, with the
following properties:

\begin{enumerate}
\def\labelenumi{\arabic{enumi}.}
\tightlist
\item
  Each dimension of the matrix represents a different feature of the
  stimuli.
\item
  The length of each dimension represents how many different values that
  feature can have.
\item
  A single item can then be represented by a binary matrix with 1s in
  the cells indicating which feature values it has and 0s everywhere
  else.
\end{enumerate}

For example, in Robert Nosofsky's (1986) test of the GCM\footnote{Nosofsky,
  R. M. (1986). Attention, similarity, and the
  identification-categorization relationship. Journal of Experimental
  Psychology: General, 115(1), 39-57.}, the stimuli were semicircles
with lines through them. The two feature dimensions were (1) circle size
and (2) line orientation. Each feature had four possible values.

Using our binary matrix representation, here is one possible stimulus
from the Nosofsky experiment:

\begin{Shaded}
\begin{Highlighting}[]
\ImportTok{import}\NormalTok{ numpy }\ImportTok{as}\NormalTok{ np}

\NormalTok{stimulus }\OperatorTok{=}\NormalTok{ np.array([[}\DecValTok{0}\NormalTok{, }\DecValTok{0}\NormalTok{, }\DecValTok{0}\NormalTok{, }\DecValTok{0}\NormalTok{],}
\NormalTok{                     [}\DecValTok{0}\NormalTok{, }\DecValTok{0}\NormalTok{, }\DecValTok{0}\NormalTok{, }\DecValTok{0}\NormalTok{],}
\NormalTok{                     [}\DecValTok{0}\NormalTok{, }\DecValTok{0}\NormalTok{, }\DecValTok{1}\NormalTok{, }\DecValTok{0}\NormalTok{],}
\NormalTok{                     [}\DecValTok{0}\NormalTok{, }\DecValTok{0}\NormalTok{, }\DecValTok{0}\NormalTok{, }\DecValTok{0}\NormalTok{]], dtype}\OperatorTok{=}\BuiltInTok{int}\NormalTok{)}
\BuiltInTok{print}\NormalTok{(stimulus)}
\end{Highlighting}
\end{Shaded}

\begin{verbatim}
[[0 0 0 0]
 [0 0 0 0]
 [0 0 1 0]
 [0 0 0 0]]
\end{verbatim}

The Nosofksy stimuli are pretty simple: they can have exactly one value
on each dimension, so this representation will always only have a single
1 in the matrix.

If we wanted to know how many of each stimulus to present our model, it
would be pretty cumbersome to maintain a long list of arrays like this
one. So we could instead maintain a single matrix that stores counts of
stimuli. That is, the values at each element will represent the number
of stimuli observed with those feature values.

For example, let's consider the ``dimensional'' condition of Nosofsky's
experiment, in which subjects saw only the stimuli on the ``diagonals''
of the matrix. Let's assume they saw each one 100 times (actually, they
did a mind-numbing 1200 trials after 2600 practice trials 🥴). We could
represent that like so:

\begin{Shaded}
\begin{Highlighting}[]
\NormalTok{training\_set }\OperatorTok{=}\NormalTok{ np.array([[}\DecValTok{100}\NormalTok{, }\DecValTok{0}\NormalTok{, }\DecValTok{0}\NormalTok{, }\DecValTok{100}\NormalTok{],}
\NormalTok{                         [}\DecValTok{0}\NormalTok{, }\DecValTok{100}\NormalTok{, }\DecValTok{100}\NormalTok{, }\DecValTok{0}\NormalTok{],}
\NormalTok{                         [}\DecValTok{0}\NormalTok{, }\DecValTok{100}\NormalTok{, }\DecValTok{100}\NormalTok{, }\DecValTok{0}\NormalTok{],}
\NormalTok{                         [}\DecValTok{100}\NormalTok{, }\DecValTok{0}\NormalTok{, }\DecValTok{0}\NormalTok{, }\DecValTok{100}\NormalTok{]], dtype}\OperatorTok{=}\BuiltInTok{int}\NormalTok{)}
\BuiltInTok{print}\NormalTok{(training\_set)}
\end{Highlighting}
\end{Shaded}

\begin{verbatim}
[[100   0   0 100]
 [  0 100 100   0]
 [  0 100 100   0]
 [100   0   0 100]]
\end{verbatim}

This still isn't ideal though. Because it doesn't represent the category
labels that subjects got during their training. We don't know which
examples belonged to which categories.

Thinking ahead to how the GCM works, let's instead represent the
training examples more like a list, so that we can iterate through them:

\begin{Shaded}
\begin{Highlighting}[]
\NormalTok{tDimensional }\OperatorTok{=}\NormalTok{ np.zeros((}\DecValTok{8}\NormalTok{,}\DecValTok{3}\NormalTok{), dtype}\OperatorTok{=}\BuiltInTok{int}\NormalTok{)}
\NormalTok{tDimensional[}\DecValTok{0}\NormalTok{] }\OperatorTok{=}\NormalTok{ [}\DecValTok{1}\NormalTok{,}\DecValTok{1}\NormalTok{,}\DecValTok{1}\NormalTok{]}
\NormalTok{tDimensional[}\DecValTok{1}\NormalTok{] }\OperatorTok{=}\NormalTok{ [}\DecValTok{1}\NormalTok{,}\DecValTok{2}\NormalTok{,}\DecValTok{2}\NormalTok{]}
\NormalTok{tDimensional[}\DecValTok{2}\NormalTok{] }\OperatorTok{=}\NormalTok{ [}\DecValTok{1}\NormalTok{,}\DecValTok{3}\NormalTok{,}\DecValTok{2}\NormalTok{]}
\NormalTok{tDimensional[}\DecValTok{3}\NormalTok{] }\OperatorTok{=}\NormalTok{ [}\DecValTok{1}\NormalTok{,}\DecValTok{4}\NormalTok{,}\DecValTok{1}\NormalTok{]}
\NormalTok{tDimensional[}\DecValTok{4}\NormalTok{] }\OperatorTok{=}\NormalTok{ [}\DecValTok{2}\NormalTok{,}\DecValTok{1}\NormalTok{,}\DecValTok{4}\NormalTok{]}
\NormalTok{tDimensional[}\DecValTok{5}\NormalTok{] }\OperatorTok{=}\NormalTok{ [}\DecValTok{2}\NormalTok{,}\DecValTok{2}\NormalTok{,}\DecValTok{3}\NormalTok{]}
\NormalTok{tDimensional[}\DecValTok{6}\NormalTok{] }\OperatorTok{=}\NormalTok{ [}\DecValTok{2}\NormalTok{,}\DecValTok{3}\NormalTok{,}\DecValTok{3}\NormalTok{]}
\NormalTok{tDimensional[}\DecValTok{7}\NormalTok{] }\OperatorTok{=}\NormalTok{ [}\DecValTok{2}\NormalTok{,}\DecValTok{4}\NormalTok{,}\DecValTok{4}\NormalTok{]}

\BuiltInTok{print}\NormalTok{(tDimensional)}
\end{Highlighting}
\end{Shaded}

\begin{verbatim}
[[1 1 1]
 [1 2 2]
 [1 3 2]
 [1 4 1]
 [2 1 4]
 [2 2 3]
 [2 3 3]
 [2 4 4]]
\end{verbatim}

In this format, each row of the matrix is an array with three elements.
The first element is the category label: 1 or 2. The second two elements
are the indices into the stimulus matrix. \textbf{Note: Here I've
deviated from Python norms to be consistent with the numbering in the
table in the original Nosofsky paper, shown below. When working with the
model, remember that indexing in Python starts with 0.}

\begin{figure}[H]

{\centering \includegraphics[width=0.5\linewidth,height=\textheight,keepaspectratio]{images/04/dimensional_table.png}

}

\caption{Training examples from the dimensional condition. Image from
Nosofsky (1986).}

\end{figure}%

\section{The Generalized Context Model
(GCM)}\label{the-generalized-context-model-gcm}

Here, I will describe a special case of the model here that applies to
tasks (like the one in the paper) with only two categories. However,
this model can easily be extended to any number of categories. The model
is defined by two key equations.

The first equation defines the probability of classifying Stimulus
\(S_i\) into Category \(C_1\) (i.e., Category 1):

\begin{equation}\phantomsection\label{eq-gcm1}{
\begin{equation}
P(C_1|S_i) = \frac{b_1 \sum_{j \in C_1} \eta_{ij}}{b_1 \sum_{j \in C_1} \eta_{ij} + (1-b_1) \sum_{k \in C_2} \eta_{ik}}
\end{equation}
}\end{equation}

The sum \(\sum_{j \in C_1}\) is a sum over all stimuli \(S_j\) that
belong to Category 1. The sum \(\sum_{j \in C_1}\) is a sum over all
stimuli \(S_k\) that belong to Category 2.

Equation~\ref{eq-gcm1} has a parameter \(b_1\) which is defined as the
response bias for Category 1. \(b_1\) can range from 0 to 1 and captures
the possibility that a subject is biased to respond with one category or
another. This is similar but not identical to a prior distribution. When
\(b_1\) is small, the model is biased toward Category 2; when \(b_1\) is
large, the model is biased toward Category 1. Note that when
\(b = 0.5\), it cancels out of the equation.

We could also include a response bias for Category 2, but the model
assumes that \(\sum_i b_i = 1\). Therefore, when there are only two
categories, \(b_2 = 1 - b_1\) which is what you see in the second term
of the denominator of Equation~\ref{eq-gcm1}.

\(\eta_{ij}\) is a function that defines how similar to \(S_i\) and
\(S_j\) are. The GCM assumes that stimuli can be represented as points
in a multi-dimensional space (in this case, a two-dimensional space) and
the similarity is defined as a function of the distance between those
two points:

\[
\eta_{ij} = e^{-c^2 \left[w_1 (f_{i1} - f_{j1})^2 + (1-w_1)(f_{i2} - f_{j2})^2 \right]}
\]

where \(f_{i1}\) and \(f_{i2}\) refer to the feature values on
dimensions 1 and 2 of \(S_i\). This equation has two parameters: \(c\)
and \(w_1\). \(c\) is a scaling parameter that affects how steeply the
exponential curve is. This will allow us to account for how different
people might have different ideas of how close two stimuli have must be
to be called similar. \(w_1\) is called the attentional weight for
dimension 1. This parameter captures how much weight is placed on
dimension 1 over dimension 2. Just like \(b_1\), \(w_1\) can range from
0 to 1, and the larger it is, the more weight is placed on dimension 1.
Similarly, we could add a \(w_2\), but the attention weights are
constrained to sum to 1.

\subsection{Generating model
predictions}\label{generating-model-predictions}

In order to generate predictions from this model, we need two things.
First, we need to define what the training stimuli are, and what their
feature values are. This is what will allow us to compute the sums in
Equation~\ref{eq-gcm1}.

Next, we need to specify the values of the three parameters \(b_1\),
\(c\), and \(w_1\). Nosofsky does this by first collecting data from
people in the experiment. That is, for every stimulus \(S_i\) we can
record the proportion of times people classified it as Category 1. In
other words, we can collect empirical estimates of \(P(C_1|S_i)\) for
all values of \(S_i\).

Then Nosofsky uses a maximum likelihood procedure for fitting the model
to the data. Conceptually, the idea is to find the set of values of the
three parameters that produce model predictions that are as close as
possible to the empirical data. (This is exactly how linear regression
works, where you have some data \((\bar{x},\bar{y})\) and want to find
the best-fitting values of \(a\) and \(b\) for the function
\(y = ax + b\) that describes the relationship between \(x\) and \(y\).)

We can define model fit by mean squared error (MSE). MSE is defined as

\[
\frac{1}{n} \sum_i (y_i - x_i)^2
\]

where \(y\) are the data and \(x\) are the model predictions. A lower
MSE means a better model fit. Finding the best values of the parameters
can be achieved through an exhaustive search of all values. For example,
we might consider all possible values of each parameter in increments of
0.1. We could then write a program that performs the following basic
algorithm:

\begin{enumerate}
\def\labelenumi{\arabic{enumi}.}
\tightlist
\item
  Set minimum MSE to \(\inf\).
\item
  For all possible values of \(b_1\), \(c\), and \(w_1\):

  \begin{enumerate}
  \def\labelenumii{\arabic{enumii}.}
  \tightlist
  \item
    Generate model predictions \(P(R_1|S_i)\) for all \(S_i\).
  \item
    Compute MSE between the empirical data and model predictions
    generated in the previous step.
  \item
    If this MSE is smaller than minimum MSE, set minimum MSE to the
    current value and store the current values of \(b_1\), \(c\), and
    \(w_1\).
  \end{enumerate}
\item
  Return the stored values of \(b_1\), \(c\), and \(w_1\).
\end{enumerate}

\subsection{Cross-validation}\label{cross-validation}

To avoid over-fitting models to our collected data, best practice is to
use an approach that splits the data into training and evaluation (or
test) sets and judge performance only on the test sets.

The basic idea, known as
\href{https://en.wikipedia.org/wiki/Cross-validation_(statistics)}{cross-validation},
comes in many forms. The principle is straightforward, but the details
are outside the scope of this book. There are existing packages in
Python that perform many of the key cross-validation steps
automatically.

\subsection{Homework 3}\label{homework-3}

\href{https://colab.research.google.com/drive/140NfyG4UV1au0722KIy54-o0f62Ifp2u?usp=sharing}{\pandocbounded{\includesvg[keepaspectratio]{index_files/mediabag/colab-badge.svg}}}

In your next homework you will implement the GCM described above. In the
homework, you will have to encode training and test stimuli and
implement the equations and model-fitting algorithm above. Additionally,
to get some experience with the behavioral side of cognitive science,
you will create a working version of a category learning experiment in
order to collect some data that you can use to compare with the GCM's
predictions.

\bookmarksetup{startatroot}

\chapter{Hierarchical generalization}\label{hierarchical-generalization}

In \hyperref[bayes]{previous examples}, there was always a finite number
of hypotheses that we were making inferences about (number of black
balls, fair or trick coin, yellow or green taxi). Sometimes, we want to
consider an infinite set of hypotheses. For example, after flipping a
coin, what is the probability of that coin coming up heads? The answer
to this question could be any number in the interval {[}0,1{]}.

\section{The Beta-Binomial model 🪙}\label{beta-binomial}

\begin{figure}[H]

{\centering \includegraphics[width=0.75\linewidth,height=\textheight,keepaspectratio]{images/05/coin-in-hand.jpg}

}

\caption{Photo by ZSun Fu on Unsplash.}

\end{figure}%

We can answer this question with a model called the Beta-Binomial model,
named for the probability distributions it uses. First, let's set up the
basic assumptions of the model.

Let \(P(\text{heads}) = \theta\). We don't know what \(\theta\) is.
After observing a sequence of coin flips \(D\), we want to estimate
\(\theta\). This can be accomplished by directly applying Bayes's rule:

\[
P(\theta|D) = \frac{P(D|\theta) P(\theta)}{P(D)}
\]

The data \(D\) in this case corresponds to the number of \(k\) heads out
of \(n\) total flips. This follows a
\href{https://en.wikipedia.org/wiki/Binomial_distribution}{Binomial
distribution}, which describes the probability of getting \(k\)
successes out of \(n\) trials, when the probability of success on each
trial is \(\theta\). We will define heads as a ``success''.

\[
\begin{align}
P(D|\theta) = P(k|\theta,n) &= \text{Bin}(k; n, \theta) \\
&= \binom{n}{k} \theta^{k} (1-\theta)^{n-k}
\end{align}
\]

The notation for the \(\text{Bin}(\cdot)\) function indicates that this
is a distribution over \(k\) (number of successes) and the distribution
has the parameters \(n\) (the total number of trials) and \(\theta\)
(the probability of a success on each trial).

We can define the prior, \(P(\theta)\), however we like. Because
\(\theta\) is a random variable that can take on any value from 0 to 1,
we cannot just say \(P(\theta) = 0.5\) like we could in earlier
examples. Instead, \(P(\theta)\) must be a probability distribution that
assigns probabilities to any value from 0 to 1. If we know nothing about
\(\theta\), we could use a Uniform(\([0,1]\)) or non-informative prior
that assigns equal probability to all values of \(\theta\).

Alternatively, a convenient choice (for reasons explained below) for
\(P(\theta)\) is the Beta distribution:

\[
P(\theta) = \text{Beta}(\theta;\alpha,\beta)
\]

The \href{https://en.wikipedia.org/wiki/Beta_distribution}{Beta
distribution} has two parameters: \(\alpha > 0\) and \(\beta > 0\).
Let's create a function that will allow us to visualize the Beta
distribution.

\begin{Shaded}
\begin{Highlighting}[]
\ImportTok{from}\NormalTok{ scipy }\ImportTok{import}\NormalTok{ stats}
\ImportTok{import}\NormalTok{ numpy }\ImportTok{as}\NormalTok{ np}
\ImportTok{import}\NormalTok{ matplotlib.pyplot }\ImportTok{as}\NormalTok{ plt}
\ImportTok{import}\NormalTok{ seaborn }\ImportTok{as}\NormalTok{ sns}

\NormalTok{sns.set\_theme()}

\KeywordTok{def}\NormalTok{ plot\_beta(a, b):}
\NormalTok{  x }\OperatorTok{=}\NormalTok{ np.linspace(}\DecValTok{0}\NormalTok{,}\DecValTok{1}\NormalTok{,num}\OperatorTok{=}\DecValTok{500}\NormalTok{)}
\NormalTok{  px }\OperatorTok{=}\NormalTok{ stats.beta.pdf(x, a, b)}
  
\NormalTok{  fig, ax }\OperatorTok{=}\NormalTok{ plt.subplots()}
\NormalTok{  sns.lineplot(x }\OperatorTok{=}\NormalTok{ x, y }\OperatorTok{=}\NormalTok{ px)}
\NormalTok{  ax.set\_xlabel(}\VerbatimStringTok{r\textquotesingle{}}\DecValTok{$}\CharTok{\textbackslash{}t}\VerbatimStringTok{heta}\DecValTok{$}\VerbatimStringTok{\textquotesingle{}}\NormalTok{)}
\NormalTok{  plt.show()}
\end{Highlighting}
\end{Shaded}

\texttt{plot\_beta} takes two arguments: \texttt{a} (\(\alpha\)), and
\texttt{b}(\(\beta\)) and plots a Beta distribution with those parameter
values.

Let's see what it looks like with a few different values.

\begin{Shaded}
\begin{Highlighting}[]
\NormalTok{plot\_beta(}\DecValTok{1}\NormalTok{,}\DecValTok{1}\NormalTok{)}
\end{Highlighting}
\end{Shaded}

\pandocbounded{\includegraphics[keepaspectratio]{05-hierarchical-bayes_files/figure-pdf/cell-3-output-1.pdf}}

When \(\alpha = \beta = 1\), the Beta distribution is identical to a
Uniform(\([0,1]\)) distribution.

\begin{Shaded}
\begin{Highlighting}[]
\NormalTok{plot\_beta(}\DecValTok{3}\NormalTok{,}\DecValTok{3}\NormalTok{)}
\end{Highlighting}
\end{Shaded}

\pandocbounded{\includegraphics[keepaspectratio]{05-hierarchical-bayes_files/figure-pdf/cell-4-output-1.pdf}}

When \(\alpha\) and \(\beta\) are greater than 1 and equal, we get a
distribution with a peak around 0.5. If we had strong prior expectations
that the coin was unbiased, we could increase the parameters even more:

\begin{Shaded}
\begin{Highlighting}[]
\NormalTok{plot\_beta(}\DecValTok{50}\NormalTok{,}\DecValTok{50}\NormalTok{)}
\end{Highlighting}
\end{Shaded}

\pandocbounded{\includegraphics[keepaspectratio]{05-hierarchical-bayes_files/figure-pdf/cell-5-output-1.pdf}}

What about when \(\alpha\) and \(\beta\) are not equal?

\begin{Shaded}
\begin{Highlighting}[]
\NormalTok{plot\_beta(}\DecValTok{4}\NormalTok{,}\DecValTok{2}\NormalTok{)}
\end{Highlighting}
\end{Shaded}

\pandocbounded{\includegraphics[keepaspectratio]{05-hierarchical-bayes_files/figure-pdf/cell-6-output-1.pdf}}

This allows us to capture skewed priors, perhaps capturing a belief that
the coin has a specific bias.

Now, what if \(\alpha\) and \(\beta\) are less than 1?

\begin{Shaded}
\begin{Highlighting}[]
\NormalTok{plot\_beta(}\FloatTok{0.5}\NormalTok{,}\FloatTok{0.5}\NormalTok{)}
\end{Highlighting}
\end{Shaded}

\pandocbounded{\includegraphics[keepaspectratio]{05-hierarchical-bayes_files/figure-pdf/cell-7-output-1.pdf}}

This might capture the belief that the coin is strongly biased, but we
aren't sure in which direction.

\subsection{Conjugate distributions}\label{conjugate-distributions}

The Beta distribution is the \emph{conjugate distribution} for the
Binomial distribution. This means that when the likelihood is a Binomial
distribution and the prior is a Beta distribution, then the posterior is
also a Beta distribution. Specifically, after making these assumptions,

\[
P(\theta|D) = \text{Beta}(\theta; \alpha + k, \beta + n-k)
\]

The parameters of the posterior distribution are (1) the sum of
\(\alpha\) from the prior and the number of observed heads and (2) the
sum of \(\beta\) from the prior and the number of observed tails. This
means that the parameters \(\alpha\) and \(\beta\) of the Beta prior
have a natural interpretation as ``virtual flips''. For example, the
larger \(\alpha\) is compared to \(\beta\), the more biased toward heads
we expect \(\theta\) to be. Additionally, the larger \(\alpha\) and
\(\beta\) are, the more certain (less diffuse) the prior is.

\subsection{Parameter estimation}\label{parameter-estimation}

Because we used a conjugate distribution, we can use our same
\texttt{plot\_beta} function to generate posterior probability
distributions after some coin flips.

Suppose you start with a fairly strong belief that a coin is fair,
represented by this distribution:

\begin{Shaded}
\begin{Highlighting}[]
\NormalTok{plot\_beta(}\DecValTok{30}\NormalTok{,}\DecValTok{30}\NormalTok{)}
\end{Highlighting}
\end{Shaded}

\pandocbounded{\includegraphics[keepaspectratio]{05-hierarchical-bayes_files/figure-pdf/cell-8-output-1.pdf}}

Now, suppose you flip a coin 20 times and it comes up heads every time.
What should you think about the bias of the coin now? According to our
model:

\begin{Shaded}
\begin{Highlighting}[]
\NormalTok{plot\_beta(}\DecValTok{30}\OperatorTok{+}\DecValTok{20}\NormalTok{,}\DecValTok{30}\NormalTok{)}
\end{Highlighting}
\end{Shaded}

\pandocbounded{\includegraphics[keepaspectratio]{05-hierarchical-bayes_files/figure-pdf/cell-9-output-1.pdf}}

As you can see, this should cause you to shift your beliefs somewhat.

This wasn't totally realistic, though. If you picked a coin off the
ground, your prior beliefs about it being biased would probably look
more like this:

\begin{Shaded}
\begin{Highlighting}[]
\NormalTok{plot\_beta(}\DecValTok{2000}\NormalTok{,}\DecValTok{2000}\NormalTok{)}
\end{Highlighting}
\end{Shaded}

\pandocbounded{\includegraphics[keepaspectratio]{05-hierarchical-bayes_files/figure-pdf/cell-10-output-1.pdf}}

What happens if you \emph{now} flipped this coin 20 times and it came up
heads every time?

\begin{Shaded}
\begin{Highlighting}[]
\NormalTok{plot\_beta(}\DecValTok{2000}\OperatorTok{+}\DecValTok{20}\NormalTok{,}\DecValTok{2000}\NormalTok{)}
\end{Highlighting}
\end{Shaded}

\pandocbounded{\includegraphics[keepaspectratio]{05-hierarchical-bayes_files/figure-pdf/cell-11-output-1.pdf}}

You might be mildly surprised, but those 20 flips wouldn't be enough to
budge your estimate about the bias of the coin by much.

Finally, let's imagine a situation in which you had a weak prior belief
that a coin was biased:

\begin{Shaded}
\begin{Highlighting}[]
\NormalTok{plot\_beta(}\DecValTok{5}\NormalTok{,}\DecValTok{1}\NormalTok{)}
\end{Highlighting}
\end{Shaded}

\pandocbounded{\includegraphics[keepaspectratio]{05-hierarchical-bayes_files/figure-pdf/cell-12-output-1.pdf}}

Now you flip the coin 100 times and it comes up heads 48 times. What
should your updated beliefs be?

\begin{Shaded}
\begin{Highlighting}[]
\NormalTok{plot\_beta(}\DecValTok{5}\OperatorTok{+}\DecValTok{48}\NormalTok{,}\DecValTok{1}\OperatorTok{+}\DecValTok{52}\NormalTok{)}
\end{Highlighting}
\end{Shaded}

\pandocbounded{\includegraphics[keepaspectratio]{05-hierarchical-bayes_files/figure-pdf/cell-13-output-1.pdf}}

As you can see, the posterior distribution shows that you should think
this coin is probably fair now. This illustrates how sufficient evidence
can override prior beliefs.

\subsection{Hypothesis averaging}\label{hypothesis-averaging}

\hyperref[generalizing]{In Chapter 3}, we solved the generalization
problem by summing over all hypotheses, weighted by their posterior
probabilities. Here, we can do something similar.

Suppose we want to know the probability of the next flip coming up
heads. In other words, we want to know \(P(\text{heads}|D)\). We can do
that by averaging over all possible values of \(\theta\):

\[
P(\text{heads}|D) = \int_\theta P(\text{heads}|\theta) \cdot P(\theta|D) d\theta = \int_\theta \theta \cdot P(\theta|D) d\theta
\]

\section{Overhypotheses 🙆}\label{overhypotheses}

Now consider a slightly different situation. You flip 19 different coins
in a row, each one time, and they all come up heads. Now you pick up a
20th coin from the same bag as the previous 19 coins. What do you think
is the probability of that 20th coin coming up heads? Is it higher than
0.5?

If you answered yes, it's probably because you formed an
\emph{overhypothesis} about the bias of the coins. After flipping all
those coins, you may have concluded that this particular set of coins is
more likely than usual to be biased. As a result, your estimate about
the probability of the 20th coin coming up heads was higher than it
otherwise would be.

\subsection{The shape bias}\label{the-shape-bias}

This coins example is pretty artificial, but the concept of
overhypotheses is one that you find in language learning. A phenomenon
known as \href{https://doi.org/10.1016/0885-2014(88)90014-7}{the shape
bias} refers to the fact that even young children are more likely to
generalize a new word based on its shape rather than other properties
like color or texture.

\begin{figure}[H]

{\centering \pandocbounded{\includegraphics[keepaspectratio]{images/05/shape_bias.gif}}

}

\caption{A common task used to test for the shape bias.}

\end{figure}%

This makes sense because objects tend to have common shapes and are less
likely to have common colors or textures.

\subsection{Modeling the learning of overhypotheses through hierarchical
Bayesian
learning}\label{modeling-the-learning-of-overhypotheses-through-hierarchical-bayesian-learning}

Charles Kemp, Andy Perfors, and Josh Tenenbaum developed
\href{http://www.charleskemp.com/papers/KempPTDevSci.pdf}{a model} of
this kind of learning. They focused on bags of black and white marbles
rather than flipping coins. They imagine a problem in which you have
many bags of marbles that you draw from. After drawing from many bags,
you draw a single marble from a new bag and make a prediction about the
proportion of black and white marbles in that bag.

The details of the model are outside the scope of this book. But the
basic idea is that the model learns at two levels simultaneously. At the
higher level, the model learns the parameters \(\alpha\) and \(\beta\)
of a Beta distribution that characterizes the proportion of black and
white marbles in each bag. As we saw above, a Beta distribution can have
a peak around a particular proportion, or it can be peaked around both 0
and 1, meaning that each bag is likely to be nearly all black or all
white.

At the lower level, the model learns the specific distribution of
marbles within a bag. If you draw 20 marbles and 5 of them are black,
you may have some uncertainty about the overall proportion in the bag,
but your best estimate will be around 5/20 (or 1/4).

Where the model excels is being able to draw inferences \emph{across}
bags. If you see many bags that are full of \emph{only} black or
\emph{only} white marbles, and then you draw a single black marble out
of a new bag, you are likely to be very confident that the rest of the
marbles in that bag are black.

But if you see many bags that have mixed proportions of black and white
marbles, and then you draw a single black marble out of a new bag, you
will be far less confident about the proportion of black marbles in that
bag. A model that doesn't make inferences at multiple levels would
struggle to draw this distinction.

\bookmarksetup{startatroot}

\chapter{Sampling assumptions}\label{sampling-assumptions}

So far, we've seen where \hyperref[weak-strong-sampling]{strong
sampling} was a logical assumption to make. But it isn't always.

For example, consider \hyperref[hormones]{the hormone problem from
Chapter 3}. If you were randomly testing people's hormone levels and
every person just happened to be healthy, clearly strong sampling
wouldn't be appropriate, because we were sampling from the full
population of both healthy and unhealthy people.

This raises a question: Are people sensitive to the process by which
data are generated?

\emph{Spoiler alert: Yes.}

\section{Word learning 💬}\label{word-learning}

Suppose you see the following collection of objects on a table.

\begin{figure}[H]

{\centering \includegraphics[width=0.7\linewidth,height=\textheight,keepaspectratio]{images/06/xu_tenenbaum_objects.png}

}

\caption{From Xu \& Tenenbaum (2007).}

\end{figure}%

Now consider the following situation:

\begin{quote}
A teacher picks three blue circles from the pile and calls them
``wugs''.
\end{quote}

\textbf{What is a wug?} Are all circles wugs? Or just blue circles?

Consider a different situation:

\begin{quote}
A teacher picks one blue circle from the pile and calls it a ``wug''.
The teacher then asks a child to choose two more (the child picks two
blue circles). The teacher confirms that they are also wugs.
\end{quote}

\textbf{What is a wug?} Are all circles wugs? Or just blue circles?

In both situations, the observed data is the same: three blue circles
labeled as wugs. What's different is the process by which the data were
generated. In one case, a knowledgeable teacher picked all the wugs; in
the other, the teacher only picked one example and a (probably not
knowledgeable) child picked the others.

To use some terms we've seen before, the teacher is using strong
sampling. The child is using something closer to weak sampling.

\begin{figure}[H]

{\centering \includegraphics[width=0.65\linewidth,height=\textheight,keepaspectratio]{images/06/xu_tenenbaum_hierarchy.png}

}

\caption{The hierarchy of objects. From Xu \& Tenenbaum (2007).}

\end{figure}%

This is based on an
\href{http://web.mit.edu/cocosci/archive/Papers/devsci07_xu_tenenbaum.pdf}{actual
study} by Fei Xu and Josh Tenenbaum, who also developed a model of the
task. They asked whether children and adults would generalize the word
``wug'' to the basic-level category (circles with lines in them) or the
subordinate category (blue circles with lines in them). Here are their
results:

\begin{figure}[H]

{\centering \includegraphics[width=0.75\linewidth,height=\textheight,keepaspectratio]{images/06/xu_tenenbaum_results.png}

}

\caption{Experimental results from Xu \& Tenenbaum (2007).}

\end{figure}%

People clearly distinguished between the teacher-driven situation and
the child-driven (learner-driven) situation, and were much more likely
to generalize ``wugs'' to the subordinate category when the teacher
picked the objects.

\section{Pedagogical sampling 🧑‍🏫}\label{pedagogical-sampling}

The teacher choosing wugs could be said to be using \emph{pedagogical
sampling}. The teacher deliberately used knowledge of the wug concept to
select informative examples to help the child learn the concept as
quickly as possible.

This is an idea
\href{http://shaftolab.com/assets/papers/shaftogg14.pdf}{explored} by
Patrick Shafto, Noah Goodman, and Thomas Griffiths (as well as Elizabeth
Bonawitz and other researchers in related work). To put it in formal
terms:

\begin{equation}\phantomsection\label{eq-pedagogical}{
\begin{equation}
P_{\text{teacher}}(d|h) \propto (P_{\text{learner}}(h|d))^\alpha
\end{equation}
}\end{equation}

Here, \(\alpha\) is a parameter that controls how optimized the
teacher's choices are. As \(\alpha \rightarrow \infty\), the teacher
will choose examples \(d\) that maximize the learner's posterior
probability.

We can figure out how the learner will update their beliefs about a
concept by directly applying Bayes rule:

\begin{equation}\phantomsection\label{eq-plearner}{
\begin{equation}
  P_{\text{learner}}(h|d) = \frac{P_{\text{teacher}}(d|h)P(h)}{\sum_{h_i} P_{\text{teacher}}(d|h_i)P(h_i)}
\end{equation}
}\end{equation}

This equation makes it clear that the teacher and learner are
inextricably linked: The teacher chooses examples to maximize the
learner's understanding, and the learner updates beliefs based on
expectations of how the teacher is choosing examples. The model is
\emph{recursive}.

As the researchers explain in their paper, the two equations above are a
system of equations that can be rearranged to create the following
equation defining how the teacher should choose examples:

\begin{equation}\phantomsection\label{eq-pteacher}{
\begin{equation}
  P_{\text{teacher}}(d|h) \propto \left( \frac{P_{\text{teacher}}(d|h) P(h)}{\sum_{h_i} P_{\text{teacher}}(d|h_i) P(h_i)} \right)^\alpha
\end{equation}
}\end{equation}

How to solve this equation? In the paper, they use an iterative
algorithm, that works like this:

\begin{enumerate}
\def\labelenumi{\arabic{enumi}.}
\tightlist
\item
  Initialize \(P_{\text{teacher}}(d|h)\) using weak sampling.
\item
  Iterate through the following steps until \(P_{\text{teacher}}(d|h)\)
  stabilizes (doesn't change from one iteration to the next):

  \begin{enumerate}
  \def\labelenumii{\arabic{enumii}.}
  \tightlist
  \item
    For each possible \(h\) and \(d\), compute
    \(P_{\text{learner}}(h|d)\) using Equation~\ref{eq-plearner} and
    \(P_{\text{teacher}}(d|h)\) values from the previous iteration.
  \item
    For each possible \(d\) and \(h\), update
    \(P_{\text{teacher}}(d|h)\) using \(P_{\text{learner}}(h|d)\) values
    from the previous step, where
    \(P_{\text{teacher}}(d|h_i) = P_{\text{learner}}(h_i|d) / \sum_{d_j} P_{\text{learner}}(h_i|d_j)\)
    (Equation~\ref{eq-pedagogical}).
  \end{enumerate}
\end{enumerate}

\subsection{The rectangle game}\label{the-rectangle-game}

The researchers tested out their model in an experiment using a simple
task called the rectangle game.

In the game, concepts are rectangular boundaries in a two-dimensional
space. The teacher knows the boundary and the learner has to figure it
out from examples given by the teacher.

\begin{figure}[H]

{\centering \includegraphics[width=0.9\linewidth,height=\textheight,keepaspectratio]{images/06/shafto_etal_expt1.png}

}

\caption{The rectangle game. From Shafto et al.~(2014).}

\end{figure}%

In one version, the teacher can only provide positive examples (``Here's
an example that's inside the boundary''). In another version, the
teacher can also provide negative examples (``Here's an example that's
not inside the boundary'').

In an experiment, they asked people to play the role of the learner in
the rectangle game. There were three conditions:

\begin{enumerate}
\def\labelenumi{\arabic{enumi}.}
\tightlist
\item
  \emph{Teaching-pedagogical learning}: People first acted as teachers,
  then as learners.
\item
  \emph{Pedagogical learning}: People acted just as learners and were
  told the examples they saw were chosen by a teacher.
\item
  \emph{Non-pedagogical learning}: People acted just as learners and
  were told the examples they saw were not chosen by a teacher.
\end{enumerate}

After showing people examples, they asked them to draw a rectangle of
their best guess about what the boundary was.

One key prediction of the model is that the most informative locations
for examples under pedagogical sampling are at the corners of the
rectangles. To test whether the learners understood this, they measured
the proportion of examples that ended up in the corners of the
rectangles people drew. The results are shown below.

\begin{figure}[H]

{\centering \pandocbounded{\includegraphics[keepaspectratio]{images/06/shafto_etal_expt1_results.png}}

}

\caption{Results from Shafto et al.'s (2014) Experiment 1.}

\end{figure}%

As you can see, when people thought a knowledgeable teacher was
providing the examples (the teaching-pedagogical and pedagogical
learning conditions), they were much more likely to draw tight
boundaries around the examples, compared to when they didn't think a
teacher was providing the examples (the non-pedagogical learning
condition). This suggests that people were indeed sensitive to the
process by which the data were generated and incorporated that
understanding into their inferences.

\subsection{Why the hypothesis space
matters}\label{why-the-hypothesis-space-matters}

What happens when teachers and learners have different assumptions about
the the problem they're looking at?

In the rectangle game, it's pretty obvious what the space of hypotheses
is: it's practically pre-specified as part of the game. But in most
real-world teaching situations, teachers have much more knowledge not
shared by the learners than just the specific concept they are trying to
convey. How might this interfere with learning?

This was the question that a group of researchers led by Rosie Aboody
\href{https://compdevlab.yale.edu/docs/2023/aboody2023-cogsci.pdf}{tested}.
They applied the model we just looked at to a problem where the space of
hypotheses was more ambiguous than in the rectangle game. In a study of
220 people (20 teachers and 200 learners), only about half of the
learners managed to fully learn the concept (a rule in this case). They
used the model to better understand what was going on.

\begin{figure}[H]

{\centering \includegraphics[width=0.85\linewidth,height=\textheight,keepaspectratio]{images/06/aboody_etal_expt1.PNG}

}

\caption{From Aboody et al.~(2023).}

\end{figure}%

They input the teachers' chosen examples into the model to see if the
model successfully learned the correct rule. What they varied was the
hypothesis space of possible rules the model considered. When the
hypothesis space was small (55 or fewer total rules), the model inferred
the correct rule for 75\% of the teachers' sets of examples (by correct,
they mean placing 95\% or more posterior probability on the correct
rule). It wasn't until they increased the hypothesis space to 95 total
rules that the percentage of times the teachers' examples led the model
to correctly infer the rule dropped to 55\%.

The researchers concluded from this that the learners were probably
considering a larger, more complex set of rules than the teachers were.
This misaligned set of assumptions likely caused the low learning rate.

In the next chapter, we'll work through one more application of people's
sensitivity to different sampling assumptions: understanding the
pragmatics of speech.

\bookmarksetup{startatroot}

\chapter{Language pragmatics}\label{pragmatics}

When we talk to each other, we often rely on shared assumptions and
context to get our points across more quickly. Linguists call this
\href{https://en.wikipedia.org/wiki/Pragmatics}{pragmatics}. For
example, consider the following exchange:

\begin{quote}
A: I'm almost out of gas.

B: There's a gas station up ahead.
\end{quote}

Person A didn't explicitly ask a question, but Person B understood that
A wouldn't have said anything at all if they didn't need something or
weren't explaining something they were about to do, so B responded with
an answer to the implied question (``Where can I get gas?'').

In this chapter, we'll see how some aspects of pragmatics can be
explained by combining Bayesian inference with two new ideas:
\emph{information theory} and \emph{decision theory} (or expected
utility).

\section{Surprisal 😯}\label{surprisal}

We'll focus on one concept from information theory: surprisal. Surprisal
captures how unexpected an observation is. Intuitively, the
unexpectedness of an observation should depend how probable it is.
Namely, the more probable it is, the less unexpected it is.

We want a function \(f(p)\) that takes a probability and gives us its
unexpectedness. Here are some reasonable constraints on \(f\):

\begin{itemize}
\tightlist
\item
  When an observation has probability 0, it is infinitely unexpected:
  \(f(0) = \infty\).
\item
  When an observation has probability 1, it is not unexpected at all:
  \(f(1) = 0\).
\item
  When we have two independent observations with probabilities \(p_1\)
  and \(p_2\), such that the probability of both occurring is
  \(p_1 \cdot p_2\), we want the total unexpectedness of observing both
  to be the sum of the unexpectedness of each observation:
  \(f(p_1 \cdot p_2) = f(p_1) + f(p_2)\).
\end{itemize}

A function that satisfies these constraints is:

\[
f(p) = -log(p)
\] Let's verify.

\begin{itemize}
\tightlist
\item
  \(-log(p) \rightarrow \infty\) as \(p \rightarrow 0\) ✅
\item
  \(-log(1) = 0\) ✅
\item
  \(-log(p_1 \cdot p_2) = -(log(p_1) + log(p_2)) = (-log(p_1))+(-log(p_2))\)
  ✅
\end{itemize}

Here's a plot of this function:

\begin{Shaded}
\begin{Highlighting}[]
\ImportTok{import}\NormalTok{ numpy }\ImportTok{as}\NormalTok{ np}
\ImportTok{import}\NormalTok{ seaborn }\ImportTok{as}\NormalTok{ sns}
\ImportTok{import}\NormalTok{ matplotlib.pyplot }\ImportTok{as}\NormalTok{ plt}

\NormalTok{p }\OperatorTok{=}\NormalTok{ np.linspace(}\FloatTok{0.001}\NormalTok{,}\DecValTok{1}\NormalTok{,num}\OperatorTok{=}\DecValTok{500}\NormalTok{)}
\NormalTok{fp }\OperatorTok{=} \OperatorTok{{-}}\NormalTok{np.log(p)}

\NormalTok{sns.set\_theme()}
\NormalTok{fig, ax }\OperatorTok{=}\NormalTok{ plt.subplots()}
\NormalTok{sns.lineplot(x }\OperatorTok{=}\NormalTok{ p, y }\OperatorTok{=}\NormalTok{ fp)}
\NormalTok{ax.set\_xlabel(}\StringTok{"p"}\NormalTok{)}
\NormalTok{ax.set\_ylabel(}\StringTok{"surprisal"}\NormalTok{)}
\NormalTok{plt.show()}
\end{Highlighting}
\end{Shaded}

\pandocbounded{\includegraphics[keepaspectratio]{07-RSA_files/figure-pdf/cell-2-output-1.pdf}}

\emph{What does surprisal have to do with language?}

When you're communicating with someone, surprisal can help you decide
what information to give the other person. For example, imagine that you
observe a number between 1 and 10 and you want to share it with someone
else. You could either say:

\begin{quote}
I saw an even number.
\end{quote}

Or you could say:

\begin{quote}
I saw two.
\end{quote}

Which would you say? The second one, duh. But why?

In formal terms, you're trying to send a signal such that the number you
observed is not surprising for the other person after receiving your
signal. That is, you're trying to minimize the surprisal of
\(P(2|\text{signal})\).

\textbf{Exercise}: Compute the surprisal of each of these two signals.

\section{The Rational Speech Act model 🍪}\label{rsa-model}

Let's apply the concept of surprisal to the following example. There
were initially two cookies in a cookie jar, you may have eaten some, and
you're telling your friend how many you had. You can say:

\begin{enumerate}
\def\labelenumi{\arabic{enumi}.}
\tightlist
\item
  ``I had \emph{some} cookies.'' (You ate one or two.)
\item
  ``I had \emph{all} the cookies.'' (You ate two.)
\item
  ``I had \emph{no} cookies.'' (You ate none.)
\item
  Nothing -- stay silent. (You could have eaten no cookies, one cookie,
  or two cookies.)
\end{enumerate}

Let's define this situation using a binary matrix. Each row is a signal
(what you say) and each column is a state (how many cookies were eaten).

\begin{Shaded}
\begin{Highlighting}[]
\NormalTok{state\_matrix }\OperatorTok{=}\NormalTok{ np.array([[}\DecValTok{0}\NormalTok{,}\DecValTok{1}\NormalTok{,}\DecValTok{1}\NormalTok{],}
\NormalTok{                         [}\DecValTok{0}\NormalTok{,}\DecValTok{0}\NormalTok{,}\DecValTok{1}\NormalTok{],}
\NormalTok{                         [}\DecValTok{1}\NormalTok{,}\DecValTok{0}\NormalTok{,}\DecValTok{0}\NormalTok{],}
\NormalTok{                         [}\DecValTok{1}\NormalTok{,}\DecValTok{1}\NormalTok{,}\DecValTok{1}\NormalTok{]])}
                         
\NormalTok{signals }\OperatorTok{=}\NormalTok{ [}\StringTok{"some"}\NormalTok{, }\StringTok{"all"}\NormalTok{, }\StringTok{"no"}\NormalTok{, }\StringTok{"silent"}\NormalTok{]}
\NormalTok{states }\OperatorTok{=}\NormalTok{ [}\StringTok{"0"}\NormalTok{,}\StringTok{"1"}\NormalTok{,}\StringTok{"2"}\NormalTok{]}
\end{Highlighting}
\end{Shaded}

We'll define a function that prints out this matrix as a heat map.

\begin{Shaded}
\begin{Highlighting}[]
\KeywordTok{def}\NormalTok{ print\_state\_matrix(m, x\_labels}\OperatorTok{=}\VariableTok{None}\NormalTok{, y\_labels}\OperatorTok{=}\VariableTok{None}\NormalTok{):}
\NormalTok{  fig, ax }\OperatorTok{=}\NormalTok{ plt.subplots()}
  \ControlFlowTok{if}\NormalTok{ (x\_labels }\KeywordTok{and}\NormalTok{ y\_labels):}
\NormalTok{    sns.heatmap(m, annot}\OperatorTok{=}\VariableTok{True}\NormalTok{, xticklabels}\OperatorTok{=}\NormalTok{x\_labels, yticklabels}\OperatorTok{=}\NormalTok{y\_labels)}
  \ControlFlowTok{else}\NormalTok{:}
\NormalTok{    sns.heatmap(m, annot}\OperatorTok{=}\VariableTok{True}\NormalTok{)}
  
\NormalTok{  plt.show()}

\CommentTok{\# Print out the matrix}
\NormalTok{print\_state\_matrix(state\_matrix, states, signals)}
\end{Highlighting}
\end{Shaded}

\pandocbounded{\includegraphics[keepaspectratio]{07-RSA_files/figure-pdf/cell-4-output-1.pdf}}

Now let's calculate how a literal listener would interpret each of these
statements. Each row in the matrix above (one for each signal) will
define a distribution over the states. That is, for state \(h\) and
signal \(d\), we want to compute:

\[
P_{\text{listener}}(h|d) \propto P(d|h) P(h)
\]

The literal listener assumes that, given a state, a signal consistent
with that state is chosen at random by the speaker. This is already
captured in the matrix above: signals consistent with a state have a
value of 1 in their cells and signals inconsistent have a 0. We will
also assume that the listener has a uniform prior over states.
Therefore, we can calculate the literal listener's probabilities by just
normalizing the rows of the matrix.

We'll do this more than once, so let's write a function to do it.

\begin{Shaded}
\begin{Highlighting}[]
\KeywordTok{def}\NormalTok{ normalize\_rows(m):}
\NormalTok{  normalized\_m }\OperatorTok{=}\NormalTok{ np.zeros((m.shape))}
\NormalTok{  row\_num }\OperatorTok{=} \DecValTok{0}
  \ControlFlowTok{for}\NormalTok{ row }\KeywordTok{in}\NormalTok{ m:}
\NormalTok{    normalized\_m[row\_num,:] }\OperatorTok{=}\NormalTok{ row }\OperatorTok{/} \BuiltInTok{sum}\NormalTok{(row)}
\NormalTok{    row\_num }\OperatorTok{+=} \DecValTok{1}
    
  \ControlFlowTok{return}\NormalTok{(normalized\_m)}
\end{Highlighting}
\end{Shaded}

\begin{Shaded}
\begin{Highlighting}[]
\NormalTok{literal\_listener }\OperatorTok{=}\NormalTok{ normalize\_rows(state\_matrix)}
\NormalTok{print\_state\_matrix(literal\_listener, states, signals)}
\end{Highlighting}
\end{Shaded}

\pandocbounded{\includegraphics[keepaspectratio]{07-RSA_files/figure-pdf/cell-6-output-1.pdf}}

Now let's calculate the pragmatic speaker. The pragmatic speaker chooses
a signal given a state. Therefore, this time, the \emph{columns} of the
matrix will be probability distributions.

Unlike the literal listener, the pragmatic speaker is trying to convey
some information, so they will choose a signal with the greatest
\emph{expected utility} for the literal listener. Expected utility in
this case is defined as the expected increase in understanding about the
true state of the world. Remember, we already have a way of quantifying
this: surprisal. So we'll define utility as negative surprisal of the
state given the signal from the point of view of the literal listener:
\(-(-log(P_{\text{listener}}(h|d)))\).

\begin{Shaded}
\begin{Highlighting}[]
\CommentTok{\# First we apply a softmax decision function}
\NormalTok{alpha }\OperatorTok{=} \DecValTok{1} \CommentTok{\# this paramater controls how close to maximizing the speaker is}

\CommentTok{\# To avoid an error, let\textquotesingle{}s avoid the cells with P(0)}
\NormalTok{nonzero\_probs }\OperatorTok{=}\NormalTok{ literal\_listener }\OperatorTok{!=} \DecValTok{0}

\NormalTok{pragmatic\_speaker }\OperatorTok{=}\NormalTok{ literal\_listener}
\NormalTok{pragmatic\_speaker[nonzero\_probs] }\OperatorTok{=}\NormalTok{ np.exp(np.log(literal\_listener[nonzero\_probs])}\OperatorTok{*}\NormalTok{alpha)}

\ControlFlowTok{for}\NormalTok{ col }\KeywordTok{in} \BuiltInTok{range}\NormalTok{(pragmatic\_speaker.shape[}\DecValTok{1}\NormalTok{]):}
\NormalTok{  pragmatic\_speaker[:,col] }\OperatorTok{=}\NormalTok{ (pragmatic\_speaker[:,col] }\OperatorTok{/}
    \BuiltInTok{sum}\NormalTok{(pragmatic\_speaker[:,col]))}

\NormalTok{print\_state\_matrix(pragmatic\_speaker, states, signals)}
\end{Highlighting}
\end{Shaded}

\pandocbounded{\includegraphics[keepaspectratio]{07-RSA_files/figure-pdf/cell-7-output-1.pdf}}

In the above code, I used something known as a ``softmax'' rather than a
true maximizing function. The exponential function captures the idea
that the speaker will choose probabilistically, generally choosing
options with more utility. This is a common assumption in both
psychology and economics because we often don't have complete
information when making decisions (or when modeling other people's
decisions), so assuming a pure maximizing function isn't always the best
choice.

However, I also included a parameter \texttt{alpha} that controls how
close to maximizing the speaker is. The larger \texttt{alpha} is, the
closer the speaker will get to maximizing utility.

Now we can calculate the probabilities for the pragmatic listener. This
will be similar to the literal listener in that the pragmatic listener
receives a signal and computes the probability of each state given a
signal. The difference is that the pragmatic listener doesn't assume the
signals are chosen at random, but that the signals are chosen by the
pragmatic speaker above.

We can compute the pragmatic listener by normalizing the matrix above
along the rows.

\begin{Shaded}
\begin{Highlighting}[]
\NormalTok{pragmatic\_listener }\OperatorTok{=}\NormalTok{ normalize\_rows(pragmatic\_speaker)}
\NormalTok{print\_state\_matrix(pragmatic\_listener, states, signals)}
\end{Highlighting}
\end{Shaded}

\pandocbounded{\includegraphics[keepaspectratio]{07-RSA_files/figure-pdf/cell-8-output-1.pdf}}

The pragmatic listener correctly infers a
\href{https://en.wikipedia.org/wiki/Scalar_implicature}{scalar
implicature}, a concept from pragmatics in linguistics. Even though ``I
had some cookies'' is consistent with eating all of the cookies, a
pragmatic listener will infer that if a person says they ate \emph{some}
then they did not eat \emph{all} (otherwise they would have said so).

\section{Does this model match human behavior?
🔵}\label{does-this-model-match-human-behavior}

In a word, yes. One of the
\href{http://langcog.stanford.edu/papers_new/frank-2012-science.pdf}{first
papers} to look at this was by Michael Frank and Noah Goodman. They
applied this model to the simple task of picking out a shape from a
small set and tested the model's predictions in a behavioral experiment.

\subsection{The task}\label{the-task}

\begin{itemize}
\tightlist
\item
  \textbf{The speaker}: ``Imagine you are talking to someone and you
  want to refer to the middle object. What word would you use:
  \emph{blue} or \emph{circle}?'' 🟦 🔵 🟩
\item
  \textbf{The listener}: ``Imagine someone is talking to you and uses
  the word \emph{blue} to refer to one of these objects. Which object
  are they talking about?'' 🟦 🔵 🟩
\end{itemize}

In the task and experiment, the researchers varied the actual set of
objects in a systematic way.

\subsection{The model}\label{the-model-1}

Their model is virtually identical to the cookie jar model we just
worked through. They assume that speakers try to choose words to
maximize the listener's utility, as measured by surprisal.

A literal listener in this case would assume that any object that is
consistent with the speaker's chosen word could have been the one they
were referring to.

With these assumptions, the researchers derive the probability of the
speaker choosing each word (details in the paper) as:

\[
P(w|r_s,C) = \frac{|w|^{-1}}{\sum_{w^\prime \in W} {|w^\prime|}^{-1}}
\]

\begin{itemize}
\tightlist
\item
  \(w\): The speaker's chosen word.
\item
  \(r_s\): The object the speaker meant to refer to.
\item
  \(C\): The set of objects.
\item
  \(|w|\): The number of objects that \(w\) could apply to.
\item
  \(W\): The set of words that apply to the object that the speaker
  meant to refer to.
\end{itemize}

Regarding \(W\), imagine that the speaker wanted to refer to the blue
circle. In that case, \(W = \{ \text{blue}, \text{circle} \}\), because
either word could apply to the object.

According to this model, speakers will tend to choose words that more
uniquely identify an object in a set. For example, in the set above,
because \emph{blue} applies to two objects (blue square and blue
circle), but \emph{circle} only applies to one, \emph{circle} will get
higher probability.

\subsection{Homework 4: Implement the
model}\label{homework-4-implement-the-model}

\href{https://colab.research.google.com/drive/1qrSDiXLG4VlBJI_XajdbAFQdcwMAqD0W?usp=sharing}{\pandocbounded{\includesvg[keepaspectratio]{index_files/mediabag/colab-badge.svg}}}

I haven't provided all the details for this model because your
assignment is to finish the implementation yourself, run some
simulations, and collect a small amount of real data to compare the
model to.

\bookmarksetup{startatroot}

\chapter{Social cognition}\label{social-cognition}

Watch the video below. What do you see?

\url{https://www.youtube.com/embed/VTNmLt7QX8E}

If you're like most people, you see more than just lifeless shapes
moving around. You see a whole drama play out involving characters with
goals and emotions.

This animation, from a
\href{https://psycnet.apa.org/doi/10.2307/1416950}{1944 study by Heider
and Simmel}, is an excellent example of our capacity for \emph{social
cognition}: thinking about other people. \textbf{We are constantly
attributing goals, beliefs, desires, and feelings to others.} The fact
that we are willing to do it for simple shapes suggests we may not be
able to help ourselves from doing it.

When someone takes an action, makes a decision, or makes a facial
expression, how do we infer what they're thinking or feeling? In this
chapter, we'll look at two simple computational approaches to answering
this question.

\section{Inverse decision-making 🍑}\label{inverse-decision-making}

All social cognition problems have the same basic character: you
\emph{observe} something about other people and you want to \emph{infer}
some underlying cause of what you observed. For example, you see someone
take an action and you want to infer the goal that motivated their
action.

The basic idea behind many approaches to understanding how people reason
about other people is that people have some kind of mental model of how
other people act and they essentially run that model backward (or invert
the model) to infer the information they don't get to see.

Consider a simple decision-making situation. Rue likes peaches more than
oranges and oranges more than apples:

\begin{quote}
🍑 \textgreater{} 🍊 \textgreater{} 🍎
\end{quote}

Now you give Rue a choice between these three fruits. Which one is she
most likely to choose?

Probably the peach, right? She might choose the orange or the apple
because she happens to be in the mood for one of those, but given the
information you have, the peach is a good guess.

Now imagine you don't know Rue's fruit preferences. She's offered the
following fruits to pick from:

\begin{quote}
🍑 🍊 🍎
\end{quote}

She chooses the peach 🍑. If you had to guess her relative preferences
for the fruits, what would you think?

At the very least, you might guess she likes peaches more than oranges
and apples. But why?

You assume that people's choices are guided by \emph{preferences} and
that people's actions are basically \emph{rational}. So seeing Rue make
one choice from a set of options gives you information about what she
likes.

Can we formalize this intuition?\footnote{The following is based on
  \href{https://journals.plos.org/plosone/article?id=10.1371/journal.pone.0092160}{a
  model originally developed by Chris Lucas and his collaborators}.}

\subsection{A preference learning
model}\label{a-preference-learning-model}

Let's start with some assumptions:

\begin{itemize}
\tightlist
\item
  Each fruit provides some utility to the person who takes it (recall
  the notion of utility introduced in \hyperref[pragmatics]{Chapter 7}).
\item
  Utilities are additive, meaning two peaches provide twice as much
  utility as one peach.
\item
  People will choose options \emph{in proportion} to their utility. This
  is the same ``softmax'' assumption we used in \hyperref[rsa-model]{the
  rational speech act model}.
\end{itemize}

These assumptions combine to give us a \emph{choice model}, a model of
how people will make choices between different options. (We've limited
our focus to fruits, but the model could be applied to anything.)

\begin{equation}\phantomsection\label{eq-choicemodel}{
\begin{equation}
p(c = o_j|\mathbf{u}, \mathbf{A}) = \frac{\exp(U_j)}{\sum^n_{k=1}\exp(U_k)}
\end{equation}
}\end{equation}

In this model, \(c\) is the choice, \(o_j\) refers to option \(j\),
\(\mathbf{A}\) is a vector specifying each option and their attributes
(or features), \(\mathbf{u}\) is a vector of utilities that the
decision-maker assigns to each attribute, and \(U_j\) is the summed
total utility in each option.

\subsubsection{Inverting the model}\label{inverting-the-choice-model}

Remember that our goal is to \emph{infer other people's preferences},
not predict their choices. So, after we see someone make a choice
between some fruits, how do we infer what fruits they like best? We can
invert Equation~\ref{eq-choicemodel} using Bayes's rule:

\begin{equation}\phantomsection\label{eq-idmodel}{
\begin{equation}
p(\mathbf{u}|c,\mathbf{A}) = \frac{p(c|\mathbf{u},\mathbf{A})p(\mathbf{u})}{p(c|\mathbf{A})}
\end{equation}
}\end{equation}

Suppose Rue has a choice between the following options:

\begin{quote}
Option 1: 🍑 🍎 🍊

Option 2: 🍌
\end{quote}

Jules has a choice between the following options:

\begin{quote}
Option 1: 🍑

Option 2: 🍎 🍊 🍌
\end{quote}

Both Rue and Jules pick Option 1, which includes a peach 🍑. Based on
their choices alone, who do you think likes peaches more? Or, to put it
a different way, which person's choice would make you more confident
they like peaches?

Let's apply the inverse decision-making model to these two choices.

We'll start by encoding these choices using a binary encoding scheme in
which each element in a list represents a fruit and this element is set
to 1 if that fruit is present in the option.

\begin{Shaded}
\begin{Highlighting}[]
\ImportTok{import}\NormalTok{ numpy }\ImportTok{as}\NormalTok{ np}

\CommentTok{\# Fruit order:}
\CommentTok{\# 1. peach}
\CommentTok{\# 2. apple}
\CommentTok{\# 3. orange}
\CommentTok{\# 4. banana}
\NormalTok{choice\_rue }\OperatorTok{=}\NormalTok{ np.array([[}\DecValTok{1}\NormalTok{,}\DecValTok{1}\NormalTok{,}\DecValTok{1}\NormalTok{,}\DecValTok{0}\NormalTok{],}
\NormalTok{                       [}\DecValTok{0}\NormalTok{,}\DecValTok{0}\NormalTok{,}\DecValTok{0}\NormalTok{,}\DecValTok{1}\NormalTok{]])}
\NormalTok{choice\_jules }\OperatorTok{=}\NormalTok{ np.array([[}\DecValTok{1}\NormalTok{,}\DecValTok{0}\NormalTok{,}\DecValTok{0}\NormalTok{,}\DecValTok{0}\NormalTok{],}
\NormalTok{                         [}\DecValTok{0}\NormalTok{,}\DecValTok{1}\NormalTok{,}\DecValTok{1}\NormalTok{,}\DecValTok{1}\NormalTok{]])}
\end{Highlighting}
\end{Shaded}

Let's define a function that calculates the choice probability in
Equation~\ref{eq-choicemodel}. This is our likelihood function.

\begin{Shaded}
\begin{Highlighting}[]
\KeywordTok{def}\NormalTok{ compute\_choice\_prob(j, u, options): }
  \CommentTok{\textquotesingle{}\textquotesingle{}\textquotesingle{}}
\CommentTok{  Returns the probability of choosing option j from the set options.}
\CommentTok{     }
\CommentTok{  Args:}
\CommentTok{    j (int): the chosen option}
\CommentTok{    u (array): a vector of utilities assigned to the}
\CommentTok{      attributes in the options}
\CommentTok{    options (array): a matrix in which each row is an}
\CommentTok{      option, and each option is a binary array of the }
\CommentTok{      attributes in that option}
\CommentTok{     }
\CommentTok{  Returns:}
\CommentTok{    p (float): probability of choosing option j}
\CommentTok{  \textquotesingle{}\textquotesingle{}\textquotesingle{}}
  
  \CommentTok{\# Compute the total utility of each option}
\NormalTok{  total\_u }\OperatorTok{=}\NormalTok{ np.}\BuiltInTok{sum}\NormalTok{(u}\OperatorTok{*}\NormalTok{options, axis}\OperatorTok{=}\DecValTok{1}\NormalTok{)}
  
\NormalTok{  p }\OperatorTok{=}\NormalTok{ np.exp(total\_u[j]) }\OperatorTok{/}\NormalTok{ np.}\BuiltInTok{sum}\NormalTok{(np.exp(total\_u))}
  \ControlFlowTok{return}\NormalTok{(p)}
\end{Highlighting}
\end{Shaded}

Let's first assume they both like all fruits equally (they get a utility
of 1 from any fruit). What is the probability of Rue and Jules making
these choices?

\begin{Shaded}
\begin{Highlighting}[]
\NormalTok{utilities }\OperatorTok{=}\NormalTok{ np.array([}\DecValTok{1}\NormalTok{,}\DecValTok{1}\NormalTok{,}\DecValTok{1}\NormalTok{,}\DecValTok{1}\NormalTok{])}

\BuiltInTok{print}\NormalTok{(}\StringTok{"Rue\textquotesingle{}s choice probability: }\SpecialCharTok{\{:.3f\}}\StringTok{"}\NormalTok{.}\BuiltInTok{format}\NormalTok{(}
\NormalTok{  compute\_choice\_prob(}\DecValTok{0}\NormalTok{, utilities, choice\_rue)))}
\BuiltInTok{print}\NormalTok{(}\StringTok{"Jules\textquotesingle{}s choice probability: }\SpecialCharTok{\{:.3f\}}\StringTok{"}\NormalTok{.}\BuiltInTok{format}\NormalTok{(}
\NormalTok{  compute\_choice\_prob(}\DecValTok{0}\NormalTok{, utilities, choice\_jules)))}
\end{Highlighting}
\end{Shaded}

\begin{verbatim}
Rue's choice probability: 0.881
Jules's choice probability: 0.119
\end{verbatim}

Unsurprisingly, Rue's choice is much more probable. What if they like
peaches four times as much as other fruits?

\begin{Shaded}
\begin{Highlighting}[]
\NormalTok{utilities }\OperatorTok{=}\NormalTok{ np.array([}\DecValTok{4}\NormalTok{,}\DecValTok{1}\NormalTok{,}\DecValTok{1}\NormalTok{,}\DecValTok{1}\NormalTok{])}

\BuiltInTok{print}\NormalTok{(}\StringTok{"Rue\textquotesingle{}s choice probability: }\SpecialCharTok{\{:.3f\}}\StringTok{"}\NormalTok{.}\BuiltInTok{format}\NormalTok{(}
\NormalTok{  compute\_choice\_prob(}\DecValTok{0}\NormalTok{, utilities, choice\_rue)))}
\BuiltInTok{print}\NormalTok{(}\StringTok{"Jules\textquotesingle{}s choice probability: }\SpecialCharTok{\{:.3f\}}\StringTok{"}\NormalTok{.}\BuiltInTok{format}\NormalTok{(}
\NormalTok{  compute\_choice\_prob(}\DecValTok{0}\NormalTok{, utilities, choice\_jules)))}
\end{Highlighting}
\end{Shaded}

\begin{verbatim}
Rue's choice probability: 0.993
Jules's choice probability: 0.731
\end{verbatim}

This increases Rue's choice probability by just over 10\%, but it
increases Jules's choice probability by over 60\%.

To invert the model, we also need to specify the prior probability,
\(p(\mathbf{u})\). We'll assume that people generally like fruit, but
they are likely to disagree about how much they like different fruits,
and there's a possibility that some people will dislike like certain
fruits.

We can capture this idea by assuming that utilities for fruits are
normally distributed with a positive mean. Let's say that they have a
mean of 2 and a standard deviation of 0.5:

\begin{Shaded}
\begin{Highlighting}[]
\ImportTok{from}\NormalTok{ scipy.stats }\ImportTok{import}\NormalTok{ norm}
\ImportTok{import}\NormalTok{ matplotlib.pyplot }\ImportTok{as}\NormalTok{ plt}
\ImportTok{import}\NormalTok{ seaborn }\ImportTok{as}\NormalTok{ sns}

\NormalTok{sns.set\_theme()}

\NormalTok{u\_mean }\OperatorTok{=} \DecValTok{2}
\NormalTok{u\_sd }\OperatorTok{=} \FloatTok{0.5}

\NormalTok{u }\OperatorTok{=}\NormalTok{ np.linspace(}\OperatorTok{{-}}\DecValTok{1}\NormalTok{,}\DecValTok{5}\NormalTok{,}\DecValTok{1000}\NormalTok{)}
\NormalTok{pu }\OperatorTok{=}\NormalTok{ norm.pdf(u,u\_mean,u\_sd)}

\NormalTok{fig, ax }\OperatorTok{=}\NormalTok{ plt.subplots(}\DecValTok{1}\NormalTok{,}\DecValTok{1}\NormalTok{)}
\NormalTok{sns.lineplot(x }\OperatorTok{=}\NormalTok{ u, y }\OperatorTok{=}\NormalTok{ pu)}
\NormalTok{ax.set\_xlabel(}\StringTok{"u"}\NormalTok{)}
\NormalTok{ax.set\_title(}\StringTok{"Prior probability of u"}\NormalTok{)}

\NormalTok{plt.show()}
\end{Highlighting}
\end{Shaded}

\pandocbounded{\includegraphics[keepaspectratio]{08-social-cognition_files/figure-pdf/cell-6-output-1.pdf}}

Additionally, we'll make the simplifying assumption that utilities for
different fruits are \emph{independent}. This means that knowing someone
likes pineapples doesn't tell you anything about how much they'll like
grapefruits.

Finally, we need a way to compute \(p(c|\mathbf{A})\). In other words,
after seeing Rue make a choice, what was the probability of
\emph{anyone} with \emph{any preferences} making that choice?

\hyperref[normalization]{Previously}, we avoided computing denominators
like this by normalizing. This works when you can enumerate the full
hypothesis space. This time, the space of hypotheses (all possible
assignments of utilities to fruits) is continuous and it's not easy to
integrate over. So we need to do something else.

\subsubsection{Monte Carlo estimation}\label{monte-carlo}

We'll use an estimation method known as
\href{https://en.wikipedia.org/wiki/Monte_Carlo_method}{Monte Carlo}
that relies on random sampling.

Here's the basic idea as it applies to our problem. We want to know the
overall probability of someone in the general population making a
certain choice. We don't know what that is, but we do know the
distribution of utilities in the general population. That's the prior
\(p(\mathbf{u})\) we just defined. So we can get a close approximation
to what we want using the following algorithm:

\begin{enumerate}
\def\labelenumi{\arabic{enumi}.}
\tightlist
\item
  Draw a random sample \(\mathbf{u}\) from \(p(\mathbf{u})\).
\item
  Compute \(p(c|\mathbf{u}, \mathbf{A})\).
\item
  Repeat many times.
\item
  Calculate the mean choice probability of all trials.
\end{enumerate}

Let's write a function to do this.

\begin{Shaded}
\begin{Highlighting}[]
\KeywordTok{def}\NormalTok{ estimate\_marginal\_likelihood(j, options, n\_samples, m, sd): }
  \CommentTok{\textquotesingle{}\textquotesingle{}\textquotesingle{}}
\CommentTok{  Returns an estimate of the probability of choosing option j}
\CommentTok{     from the set of options using Monte Carlo estimation.}
\CommentTok{     }
\CommentTok{   Args:}
\CommentTok{    j (int): the chosen option}
\CommentTok{    options (array): a matrix in which each row is an}
\CommentTok{     option, and each option is a binary array of the }
\CommentTok{     attributes in that option}
\CommentTok{    n\_samples (int): number of Monte Carlo samples to}
\CommentTok{     collect}
\CommentTok{    m: mean of utility prior}
\CommentTok{    sd: std. dev. of utility prior}
\CommentTok{      }
\CommentTok{       }
\CommentTok{   Returns:}
\CommentTok{    (float): probability of choosing option j}
\CommentTok{  \textquotesingle{}\textquotesingle{}\textquotesingle{}}
  
\NormalTok{  n\_attributes }\OperatorTok{=} \BuiltInTok{len}\NormalTok{(options[}\DecValTok{0}\NormalTok{])}
\NormalTok{  mc\_samples }\OperatorTok{=}\NormalTok{ np.zeros(n\_samples)}
  
  \ControlFlowTok{for}\NormalTok{ i }\KeywordTok{in} \BuiltInTok{range}\NormalTok{(n\_samples):}
    \CommentTok{\# draw random u sample}
\NormalTok{    u }\OperatorTok{=}\NormalTok{ np.random.normal(m,sd,n\_attributes)}
\NormalTok{    mc\_samples[i] }\OperatorTok{=}\NormalTok{ compute\_choice\_prob(j, u, options)}
  
  \ControlFlowTok{return}\NormalTok{(np.mean(mc\_samples))}
\end{Highlighting}
\end{Shaded}

Let's compare the probability of someone making
\hyperref[inverting-the-choice-model]{the choices that Rue and Jules
made}.

\begin{Shaded}
\begin{Highlighting}[]
\BuiltInTok{print}\NormalTok{(}\StringTok{"Rue\textquotesingle{}s choice probability: }\SpecialCharTok{\{:.3f\}}\StringTok{"}\NormalTok{.}\BuiltInTok{format}\NormalTok{(}
\NormalTok{  estimate\_marginal\_likelihood(}\DecValTok{0}\NormalTok{, choice\_rue, }\DecValTok{2000}\NormalTok{, u\_mean, u\_sd)))}
\BuiltInTok{print}\NormalTok{(}\StringTok{"Jules\textquotesingle{}s choice probability: }\SpecialCharTok{\{:.3f\}}\StringTok{"}\NormalTok{.}\BuiltInTok{format}\NormalTok{(}
\NormalTok{  estimate\_marginal\_likelihood(}\DecValTok{0}\NormalTok{, choice\_jules, }\DecValTok{2000}\NormalTok{, u\_mean, u\_sd)))}
\end{Highlighting}
\end{Shaded}

\begin{verbatim}
Rue's choice probability: 0.972
Jules's choice probability: 0.028
\end{verbatim}

Unsurprisingly, Rue's choice has a much higher baseline probability
(because she picked three fruits over one and Jules picked one fruit
over three).

\subsubsection{Putting it all together}\label{putting-it-all-together}

Now we'll infer Rue's and Jules's preferences from their choices using
Monte Carlo estimation again.

\begin{Shaded}
\begin{Highlighting}[]
\CommentTok{\# Set random state for reproducibility}
\NormalTok{np.random.RandomState(}\DecValTok{2022}\NormalTok{) }

\NormalTok{n\_samples }\OperatorTok{=} \DecValTok{2000} \CommentTok{\# number of Monte Carlo samples}
\NormalTok{n\_attributes }\OperatorTok{=} \BuiltInTok{len}\NormalTok{(choice\_rue[}\DecValTok{0}\NormalTok{])}
\NormalTok{peach }\OperatorTok{=} \DecValTok{0}
\NormalTok{banana }\OperatorTok{=} \DecValTok{3}

\NormalTok{marginal\_likelihood\_rue }\OperatorTok{=}\NormalTok{ estimate\_marginal\_likelihood(}
  \DecValTok{0}\NormalTok{, choice\_rue, n\_samples, u\_mean, u\_sd)}
\NormalTok{marginal\_likelihood\_jules }\OperatorTok{=}\NormalTok{ estimate\_marginal\_likelihood(}
  \DecValTok{0}\NormalTok{, choice\_jules, n\_samples, u\_mean, u\_sd)}

\NormalTok{mc\_samples\_rue }\OperatorTok{=}\NormalTok{ np.zeros(n\_samples)}
\NormalTok{mc\_samples\_jules }\OperatorTok{=}\NormalTok{ np.zeros(n\_samples)}

\ControlFlowTok{for}\NormalTok{ i }\KeywordTok{in} \BuiltInTok{range}\NormalTok{(n\_samples):}
  \CommentTok{\# Draw u samples for Rue and Jules}
\NormalTok{  u\_rue }\OperatorTok{=}\NormalTok{ np.random.normal(u\_mean,u\_sd,n\_attributes)}
\NormalTok{  u\_jules }\OperatorTok{=}\NormalTok{ np.random.normal(u\_mean,u\_sd,n\_attributes)}
  
  \CommentTok{\# Compute posterior for peach utility for each person}
\NormalTok{  pu\_rue }\OperatorTok{=}\NormalTok{ ((compute\_choice\_prob(}\DecValTok{0}\NormalTok{, u\_rue, choice\_rue) }\OperatorTok{*}\NormalTok{ u\_rue) }\OperatorTok{/} 
\NormalTok{    marginal\_likelihood\_rue)}
\NormalTok{  pu\_jules }\OperatorTok{=}\NormalTok{ ((compute\_choice\_prob(}\DecValTok{0}\NormalTok{, u\_jules, choice\_jules) }\OperatorTok{*}\NormalTok{ u\_jules) }\OperatorTok{/}
\NormalTok{    marginal\_likelihood\_jules)}

\NormalTok{  mc\_samples\_rue[i] }\OperatorTok{=}\NormalTok{ pu\_rue[peach]}
\NormalTok{  mc\_samples\_jules[i] }\OperatorTok{=}\NormalTok{ pu\_jules[peach]}
\end{Highlighting}
\end{Shaded}

Let's look at the results.

\begin{Shaded}
\begin{Highlighting}[]
\ImportTok{import}\NormalTok{ pandas }\ImportTok{as}\NormalTok{ pd}

\NormalTok{samples }\OperatorTok{=}\NormalTok{ pd.DataFrame(data }\OperatorTok{=}\NormalTok{ \{}\StringTok{\textquotesingle{}Rue\textquotesingle{}}\NormalTok{: mc\_samples\_rue,}
                               \StringTok{\textquotesingle{}Jules\textquotesingle{}}\NormalTok{: mc\_samples\_jules\})}
\CommentTok{\# reorganize the data}
\NormalTok{samples }\OperatorTok{=}\NormalTok{ samples.melt(value\_vars}\OperatorTok{=}\NormalTok{[}\StringTok{"Jules"}\NormalTok{,}\StringTok{"Rue"}\NormalTok{], }
\NormalTok{  var\_name}\OperatorTok{=}\StringTok{"person"}\NormalTok{, value\_name}\OperatorTok{=}\StringTok{"u"}\NormalTok{)}
\CommentTok{\# filter out extreme samples}
\NormalTok{samples\_filtered }\OperatorTok{=}\NormalTok{ samples[samples[}\StringTok{"u"}\NormalTok{] }\OperatorTok{\textless{}} \DecValTok{50}\NormalTok{]}

\NormalTok{sns.displot(data}\OperatorTok{=}\NormalTok{samples\_filtered, x}\OperatorTok{=}\StringTok{"u"}\NormalTok{, hue}\OperatorTok{=}\StringTok{"person"}\NormalTok{, kind}\OperatorTok{=}\StringTok{"kde"}\NormalTok{)}

\NormalTok{plt.show()}
\end{Highlighting}
\end{Shaded}

\pandocbounded{\includegraphics[keepaspectratio]{08-social-cognition_files/figure-pdf/cell-10-output-1.pdf}}

The model's inferences here require some interpretation. The model is
more certain about Rue having a preference for peach, but the
distribution on the Jules's preference for peach has a much longer tail.

This makes some sense. Because of our mostly positive prior
distribution, it's reasonable to assume by default that Rue likes
peaches.

Jules's choice is a bit odd, so it's harder to make sense of. One way to
explain it is that she simply made a mistake (after all,
\hyperref[monte-carlo]{as we saw above}, it's an improbable choice).
Another way to explain it is by assuming she has a very strong
preference for peaches. (Of course, the model doesn't take into account
the countless other reasons Jules might prefer one fruit over three.)

\subsection{Homework 5: Your turn}\label{homework-5-your-turn}

\href{https://colab.research.google.com/drive/1hGTsIzEeT5YtS1X2aA625YZDI8OSLCSU?usp=sharing}{\pandocbounded{\includesvg[keepaspectratio]{index_files/mediabag/colab-badge.svg}}}

In your next assignment, you'll fully implement a version of this model
and compare it to data I collected (along with Chris Lucas and Charles
Kemp).

\begin{figure}[H]

{\centering \includegraphics[width=0.75\linewidth,height=\textheight,keepaspectratio]{images/08/cards.jpg}

}

\caption{Cards used in the preference learning task from Jern, Lucas, \&
Kemp (2017).}

\end{figure}%

\section{Naïve utility calculus 🧮}\label{nauxefve-utility-calculus}

Suppose Rue and Jules both make a choice between the same two fruits:

\begin{quote}
Option 1: 🍑

Option 2: 🍌
\end{quote}

They both choose the peach 🍑. The difference is that Rue and Jules had
to exert a different degree of effort to make their choices. For Rue,
both fruits were placed in a bowl in front of her. For Jules, the peach
was in a bowl in front her and the banana was on a high shelf in another
room.

In both cases, you'd probably assume they like peaches, but you might
conclude there's weaker evidence that Jules likes them. She may have
chosen the peach because it was just convenient.

This (pretty contrived) example shows how people's choices are a
function of both \emph{rewards} (or preferences) and \emph{costs}. And
when we're thinking about other people's behavior, we take their
potential rewards and costs into account to understand why they're doing
things.

This idea was best articulated and formalized by
\href{https://compdevlab.yale.edu/docs/Jara-EttingerGweonShulzTenenbaum_TiCS.pdf}{Julian
Jara-Ettinger and his collaborators}. And they incorporated it into a
computational model.

Specifically, suppose a person has a plan \(p\) to achieve an outcome
\(o\). We can define their utility \(U\) as:

\[
U(o,p) = R(o) - C(p)
\]

where \(R(\cdot)\) is their reward from the outcome and \(C(\cdot)\) is
the cost they incur from carrying out the plan.

\subsection{Markov decision processes}\label{markov-decision-processes}

Earlier, we considered single decisions that happened in isolation. But
lots of behavior involves a series of small decisions that happen in a
sequence: the order in which someone prepares and cooks a dish, the
route someone takes to get from their workplace to their home, which
sections of a book chapter they write first.

A series of decisions is a plan \(p\). And a useful computational
framework for choosing optimal plans is a
\href{https://en.wikipedia.org/wiki/Markov_decision_process}{Markov
decision process}.

To be concrete, let's focus on the specific environment in
\href{https://compdevlab.yale.edu/docs/2020/cogpsych_NUC.pdf}{Jara-Ettinger
et al (2020)}: a 7x7 grid in which agents can move in any direction, one
step at a time.

\begin{figure}[H]

{\centering \includegraphics[width=0.6\linewidth,height=\textheight,keepaspectratio]{images/08/jara-ettinger-grid.png}

}

\caption{Example of a grid and an agent path from Jara-Ettinger et al
(2020).}

\end{figure}%

The locations in the grid are states \(S\). Whenever an agent is in a
state \(s\), it can take an action \(a \in A\), moving to any adjacent
state. The agent has some goal state in mind (in the figure above, it's
the green square). What we want is an optimal \emph{policy} that takes a
state and returns an optimal action for getting toward the goal state.
This can be computed as follows.

First, we compute each state's optimal value \(V^*(s)\), using this
recursive equation:

\[
V^*(s) = \text{max}_a\gamma \sum_{s^\prime} P_{s,a}(s^\prime) V^*(s^\prime) + R(a,s) - C(a,s)
\]

where \(P_{s,a}(s^\prime)\) is ``the probability that the agent will be
in state \(s^\prime\) when she takes action \(a\) in state \(s\) and
\(\gamma \in {0,1}\).''

Essentially, this equation requires going through every state in the
grid, summing up its rewards and costs, plus the expected downstream
rewards and costs of that state.

Solving this recursive function is outside the scope of this book, but
there are standard methods for doing it and even
\href{https://pymdptoolbox.readthedocs.io/en/latest/index.html}{Python
packages} for solving MDPs.

The action policy is then defined as:

\begin{equation}\phantomsection\label{eq-optimalpolicy}{
\begin{equation}
p(a|s) \propto \exp \left( \sum_{s^\prime} P_{s,a}(s^\prime) V^*(s^\prime) \right)
\end{equation}
}\end{equation}

In their model, they once again introduce a notion of a ``softmax''
choice rule, rather than assume agents will be perfectly optimal. This
makes sense because the goal is to model how people reason about others,
and people often have incomplete information about other people.

\subsection{Inverting the MDP}\label{inverting-the-mdp}

After seeing an agent take a series of actions \(\mathbf{a}\), we can
infer their costs and rewards pretty much the same way we did before, by
applying Bayes's rule.

\[
p(C,R|\mathbf{a}) \propto p(\mathbf{a}|C,R) p(C,R)
\]

In the model, they use uniform priors over costs and rewards. Computing
\(p(\mathbf{a}|C,R)\) means just applying
Equation~\ref{eq-optimalpolicy} repeatedly for each action the agent
took in their sequence. (The researchers also modeled a notion of
sub-goals that I've omitted.)

\subsection{Results}\label{results-1}

The researchers tested their model in a series of experiments in which
people saw paths that agents took across the grid and then had to rate
the costs and rewards the agent assigned to different states.

Results from one experiment are below.

\begin{figure}[H]

{\centering \includegraphics[width=0.9\linewidth,height=\textheight,keepaspectratio]{images/08/jara-ettinger-results.png}

}

\caption{Partial results from Experiment 1b of Jara-Ettinger et al
(2020).}

\end{figure}%

The cost and reward ratings were separately normalized so that they had
a mean of zero. Then they were combined in the plots.

Compared to simpler alternative models, the full naïve utility calculus
model provided the best fit to people's judgments.

\bookmarksetup{startatroot}

\chapter{Iterated learning}\label{iterated-learning}

If someone told you they defecated their water before drinking it, you
probably wouldn't be too impressed. Unless you lived in the 1600s, when
``defecate''
\href{https://www.mentalfloss.com/article/54770/15-words-dont-mean-what-they-used}{meant
``to purify something''}.

Language changes. Words change. Concepts change. Is there any way to
predict the conceptual drift that is always happening in our culture?

Yeah, maybe, kinda! But first, some math.

\section{Markov chains 🔗}\label{markov-chains}

Words and concepts change over \emph{time}. \textbf{Markov chains}
provide a useful modeling framework for time-dependent data.

As the name suggests, a Markov chain is a chain -- specifically of
states. At each time step, it moves to a new state. A Markov chain meets
the following conditions:

\begin{enumerate}
\def\labelenumi{\arabic{enumi}.}
\tightlist
\item
  A system can be in one of a finite number of states.
\item
  The state at each step of the chain depends only on the previous
  state. (This dependence can be probabilistic.)
\end{enumerate}

For example, consider the game Telephone, where one person whispers a
word to the next person, they whisper what they hear to the next person,
and so on. The set of possible words are the states (finite but large).

The word that Person \(n+1\) hears only depends on the word that Person
\(n\) whispered. But there might be some miscommunication: that is, the
probability that Person \(n+1\) accurately hears the word could be less
than 1, with similar-sounding words getting higher probability than
dissimilar words.

\subsection{Examples}\label{examples}

Consider a very simple chain of colors that are either red or blue.
Here, the state space \(S = \{ \text{red, blue} \}\). To specify this
Markov chain, we need to define the \emph{transition probabilities}: the
probabilities of moving from each state to every other state.

Let \(t_{ij}\) be the probability of transitioning from state \(i\) to
state \(j\). For example, if \(t_{RB} = 0.3\) then the probability of
transitioning from red to blue is 0.3. We can then write the transition
probabilities in a matrix.

\subsubsection{Example 1}\label{markov-example1}

Here's an example transition probability matrix for the red/blue state
space.

\begin{longtable}[]{@{}lll@{}}
\toprule\noalign{}
& R & B \\
\midrule\noalign{}
\endhead
\bottomrule\noalign{}
\endlastfoot
R & 0 & 1 \\
B & 1 & 0 \\
\end{longtable}

In this example, \(t_{RB} = t_{BR} = 1\). This means that the chain will
deterministically alternate between red and blue. Note that the rows of
the transition matrix must sum to 1 because they represent all the
possible states that the chain might transition to given a current
state. Although the columns also sum to 1 in this example, this is not a
requirement.

\emph{Example sequence}: {R} {B} {R} {B} {R} {B} \ldots{}

\subsubsection{Example 2}\label{markov-example2}

\begin{longtable}[]{@{}lll@{}}
\toprule\noalign{}
& R & B \\
\midrule\noalign{}
\endhead
\bottomrule\noalign{}
\endlastfoot
R & 0.25 & 0.75 \\
B & 0.75 & 0.25 \\
\end{longtable}

In this example, we'd expect to the sequence to flip back and forth
between red and blue, but not deterministically. At every step, the
sequence is more likely to flip than to stay in the current color.

\emph{Example sequence}: {R} {B} {R} {R} {B} {R} {R} {B} {R} {B} {R}
\ldots{}

\subsubsection{Example 3}\label{markov-example3}

\begin{longtable}[]{@{}lll@{}}
\toprule\noalign{}
& R & B \\
\midrule\noalign{}
\endhead
\bottomrule\noalign{}
\endlastfoot
R & 0.75 & 0.25 \\
B & 0.75 & 0.25 \\
\end{longtable}

In this example, no matter what state we're in, there is a 0.75
probability of moving to red and a 0.25 probability of moving to blue.
Therefore, in this example, we'd expect the chain to spend about 3/4 of
its time in the red state and 1/4 of its time in the blue state. This
example makes clear that the columns do not need to sum to 1.

\emph{Example sequence}: {R} {R} {B} {B} {R} {R} {B} {B} {R} {R} {R}
\ldots{}

\subsection{Stationary distributions}\label{stationary-distributions}

A common question for a given Markov chain is what is its so-called
\emph{stationary distribution}, the proportion of time the chain will
spend in each state if it runs for a very long time. For Markov chains
with a small number of states, computing the stationary distribution is
easy to do with some algebra. Consider the following example.

Let's define a transition matrix for Markov chain with two states,
\(s_1\) and \(s_2\):

\[
\begin{equation}
T = \left(
\begin{matrix}
t_{11} & t_{12} \\
t_{21} & t_{22}
\end{matrix} \right)
\end{equation}
\]

The stationary distribution corresponds to the probability of the chain
being in state \(s_1\) versus \(s_2\). We will define \(\theta_1\) as
the probability of the chain being in state \(s_1\) and \(\theta_2\) as
the probability of the chain being in state \(s_2\). We can define
\(\theta_1\) and \(\theta_2\) recursively as follows:

\[
\begin{eqnarray}
\theta_1 = t_{11} \theta_1 + t_{21} \theta_2 \\
\theta_2 = t_{22} \theta_2 + t_{12} \theta_1 \\
\end{eqnarray}
\]

These equations essentially say that the probability of being in a state
is the probability of being in that state and staying in that state plus
the probability of being in the other state and switching states. We
will then use the following facts to solve for the \(\theta\)s:

\[
\begin{eqnarray}
\theta_1 + \theta_2 = 1 \\
t_{11} + t_{12} = 1 \\
t_{21} + t_{22} = 1
\end{eqnarray}
\]

Using some algebra, we get the following expressions:

\[
\begin{eqnarray}
\theta_1 = \frac{t_{21}}{t_{12}+t_{21}} \\
\theta_2 = \frac{t_{12}}{t_{21}+t_{12}}
\end{eqnarray}
\]

We can see that these expressions give sensible results by applying them
to the red-blue examples above.

In \hyperref[markov-example1]{Example 1}, unsurprisingly,
\(\theta_R = \theta_B = 0.5\). Because the sequence alternates, it will
be red half the time and blue half the time. In
\hyperref[markov-example2]{Example 2}, again
\(\theta_R = \theta_B = 0.5\). Even though the alternating is not
deterministic, no matter what state the sequence is in, there is an
equal probability that it will switch to the other state. In
\hyperref[markov-example3]{Example 3}, \(\theta_R = 0.75\) and
\(\theta_B = 0.25\).

\section{Hidden Markov models 🙈}\label{hmms}

Now let's consider a more complicated situation in which each step in a
Markov chain generates an observation. A learner gets to see the
observations, but not the states themselves. In this situation, we say
that the states are ``hidden'', hence the name \textbf{hidden Markov
model}.

Hidden Markov models are sometimes used to model language production,
where the hidden states are the parts of speech (like noun, verb,
adjective, or more complex things like noun phrases, objects, and
subjects). What we actually get to observe as readers and listeners,
however, are words (like ``dog'', ``eats'', and ``play''). In English,
many words are ambiguous and can potentially belong to different parts
of speech. ``Play'', for example could be a noun (``The actors performed
in a play'') or a verb (``The children play in the sand''). We are
constantly decoding the language we observe to infer the parts of speech
of the words.

\subsection{Formal definition}\label{formal-definition}

A hidden Markov model includes the following pieces:

\begin{enumerate}
\def\labelenumi{\arabic{enumi}.}
\tightlist
\item
  A set of states \(S = \{s_1, \ldots, s_N\}\)
\item
  A set of observations \(O = \{o_1, \ldots, o_M\}\)
\item
  A set of initial state probabilities
  \(\Pi = \{\pi_1, \ldots, \pi_N\}\), where \(\pi_i\) is the probability
  of state \(i\) begin the first state in the sequence
\item
  A matrix of state transition probabilities \(T\), where \(t_{ij}\) in
  the matrix is the probability of transitioning from state \(s_i\) to
  state \(s_j\)
\item
  A matrix of ``emission'' probabilities \(B\), where \(b_{ik}\) in the
  matrix is the probability of producing observation \(o_k\) while in
  state \(s_i\)
\item
  An observation sequence \(Y = (y_1, \ldots, y_T)\), where \(T\) is the
  length of the observed sequence of data
\item
  An state sequence that is not observed \(X = (x_1, \ldots, x_T)\)
\end{enumerate}

\begin{figure}[H]

{\centering \includegraphics[width=0.85\linewidth,height=\textheight,keepaspectratio]{images/09/hmm.png}

}

\caption{A hidden Markov model. Source: Wikipedia (public domain).}

\end{figure}%

\subsection{Inference}\label{inference}

Usually the purpose of using an HMM is to infer \(X\) from \(Y\). There
is a standard algorithm for doing this called
\href{https://en.wikipedia.org/wiki/Viterbi_algorithm}{the Viterbi
algorithm}, the details of which are beyond the scope of this book.

There are many existing programming packages available for working with
HMMs. On Python package is
\href{https://github.com/hmmlearn/hmmlearn}{hmmlearn}, which includes a
\href{https://hmmlearn.readthedocs.io/en/stable/api.html\#hmmlearn.base.BaseHMM.predict}{\texttt{predict()}}
function, which will predict the most probable sequence of hidden states
\(X\), given a sequence of observations \(Y\) and a fully specified
Hidden Markov model. Note that this function will not always be correct.
If the model is probabilistic and there are some observations that can
be produced in more than one state (for example, the word ``play'' could
be produced by either the noun or verb state), there will be some cases
where it is impossible to be certain which state an observation came
from.

\section{Cultural transmission 🗣}\label{cultural-transmission}

So what does all this have to do with ``defecate'' changing meaning over
time?

Well, you can think of communication over time as one giant long-term
game of Telephone. And an important question you can then ask is: What
is the stationary distribution of these chains?

This is the subject of a line research in computational cognitive
science. The basic model combines some assumptions we've seen before
with some assumptions of hidden Markov models you just learned about in
this chapter.

First, assume that people are trying to learn some concept by seeing an
example \(x_i\) and getting a label for it \(y_i\) and then trying to
generalize that label to other instances by forming a hypothesis \(h\)
about what other things \(y_i\) applies to. This is the problem we
encountered in the \hyperref[generalization]{generalization chapter},
and we can assume that each person forms their beliefs about \(h\) using
Bayesian inference.

Now let's add in the communication factor. Each person \(n+1\) gets a
label \(y_n\) from someone else Person \(n\), forms their own hypothesis
\(h_{n+1}\) then labels a new object \(x_{n+1}\) for the next person.
This process repeats, creating a Markov chain.

To simplify this chain, what we end up with is a sequence of states
where each state is a hypothesis about the concept.

\emph{What will this chain converge to?}

\textbf{It turns out that the stationary distribution of this chain is
the prior} \(P(h)\). Because people make mistakes and communication is
noisy, over time, the labels \(y_n\) become virtually useless and people
will converge to their shared prior beliefs.

\subsection{Using iterated learning to identify people's
priors}\label{using-iterated-learning-to-identify-peoples-priors}

For this reason, we can use a kind of Telephone game as a way to find
out what people's shared prior beliefs are.

\href{https://link.springer.com/content/pdf/10.3758/BF03194066.pdf}{A
study lead by Michael Kalish} tested this idea for function learning:
learning to predict \(y\) values from \(x\) values.

In the task, people were given \(x\) values and had to predict the
corresponding \(y\) values using a slider. They got feedback on 50
training examples. Then they did an additional 25 trials without
feedback.

Those 25 test trials were doubled and used as the 50 training examples
for the next person, and so on. What functions would people eventually
learn?

The figure below shows several representative chains. The researchers
found that regardless of what the initial training data was, the chains
usually converged to a linear increasing function, suggesting that this
is what most people were expecting.

\begin{figure}[H]

{\centering \pandocbounded{\includegraphics[keepaspectratio]{images/09/kalishetal_functionlearning_results.png}}

}

\caption{Iterated function learning results from Kalish et al.~(2007).}

\end{figure}%

This same method has been applied in other domains as well. For example,
in \href{https://suchow.io/assets/docs/suchow2016dzp.pdf}{a study lead
by Jordan Suchow}, people had to reproduce as accurately as possible the
positions of characters in a word. The initial positions were randomly
positioned, and what each person reproduced would be passed to the next
person in the chain.

\begin{figure}[H]

{\centering \pandocbounded{\includegraphics[keepaspectratio]{images/09/suchowetal_letterspacing.png}}

}

\caption{Iterated typesetting results from Suchow et al.~(2016).}

\end{figure}%

The researchers found that, over time, people drifted much closer to
equal spaced letters, which was much more legible. Moreover, if they
averaged the responses across all chains (DZP2 in the figure), they
found that people's responses were even better than equal spacing,
getting closer to Linotype, recommended by designers.

\bookmarksetup{startatroot}

\chapter{Causal inference}\label{causal-inference}

The world is full of weird and interesting relationships. Like this one.

\begin{figure}[H]

{\centering \includegraphics[width=0.8\linewidth,height=\textheight,keepaspectratio]{images/10/pikachu-puerto-rico.png}

}

\caption{Correlation between ``surprised Pikachu'' meme and number of
middle school teachers in Puerto Rico. Source: tylervigen.com.}

\end{figure}%

Which of these weird relationships are \emph{causal} relationships and
how can we tell? Here's an AI-generated explanation for why the rise in
the meme's popularity might be responsible for the uptick in teachers in
Puerto Rico:

\begin{quote}
As the `surprised Pikachu' meme gained traction, it led to an uptick in
internet usage. This surge in online activity put a higher demand on the
telecommunication infrastructure in Puerto Rico. To meet this demand,
more tech companies invested in the region, creating job opportunities.
With the expanding job market, there was a need for additional educators
to support the growing number of families relocating to Puerto Rico. As
a result, the popularity of the meme indirectly contributed to an
increase in the number of middle school teachers in Puerto Rico.
\end{quote}

\includegraphics[width=0.4\linewidth,height=\textheight,keepaspectratio]{images/10/surprised pikachu.jpg}

Causal inference is, in fact, a major area of research in statistics and
machine learning. But we'll just focus on the question of how people
decide what causes what.

We'll use a computational framework for making optimal probabilistic
causal judgments called Bayesian networks, or \textbf{Bayes nets} for
short. Comparing people's judgments to Bayes net predictions allows us
to see how optimal (or not) people are. Additionally, as we'll see,
Bayes nets are well suited for \emph{intervention}, which is one way
that people learn about causes.

\section{Bayes nets ➡️}\label{bayes-nets}

A Bayes net is a graph that describes the dependencies between all the
variables in a situation.

For example, let's make a Bayes net for the problem in Chapter 5 of
\hyperref[beta-binomial]{inferring the bias of a coin}. In that problem,
there were three key variables: the bias \(\theta\), the total number of
flips \(n\), and the number of heads \(k\). We can represent this as a
Bayes net.

\begin{figure}[H]

{\centering \pandocbounded{\includegraphics[keepaspectratio]{images/10/beta_binomial_bayesnet1.png}}

}

\caption{Bayes net representation of the coin bias problem.}

\end{figure}%

In Bayes nets, shaded nodes represent variables that are known or
observed, and unshaded nodes represent variables that are unknown. We
know how many times the coin was flipped and came up heads, but we don't
directly know the bias \(\theta\).

We could extend this Bayes net to capture the \emph{generalization}
problem of predicting the outcome of the next coin flip \(x\).

\begin{figure}[H]

{\centering \pandocbounded{\includegraphics[keepaspectratio]{images/10/beta_binomial_bayesnet2.png}}

}

\caption{Bayes net representation of the biased coin generalization
problem.}

\end{figure}%

A complete Bayes net also specifies a probability distribution for each
variable. In the example above, \(k\) is a function of \(\theta\) and
\(n\). As we learned in \hyperref[beta-binomial]{Chapter 5}, this is a
Binomial distribution. We also need to specify a prior probability over
any unknown variables, like \(\theta\). Previously we assumed it was
distributed according to a Beta distribution. \(x\) is just a single
coin flip -- it's a special case of a Binomial distribution called a
Bernoulli distribution in which \(n = 1\). To sum up:

\begin{itemize}
\tightlist
\item
  \(\theta \sim \text{Beta}(\alpha, \beta)\)
\item
  \(k \sim \text{Binomial}(n,\theta)\)
\item
  \(x \sim \text{Binomial}(n=1, \theta)\)
\end{itemize}

\section{Causal intervention 🪚}\label{causal-intervention}

Consider the following two Bayes nets.

\begin{figure}[H]

{\centering \pandocbounded{\includegraphics[keepaspectratio]{images/10/three_var_nets.png}}

}

\caption{Two Bayes nets with three variables.}

\end{figure}%

To be concrete, let's say that variables \(s\) and \(x\) represent
levels of hormones sonin and xanthan, respectively. Variable \(z\) is an
unknown variable.

The \emph{common cause} network is so-called because sonin (\(s\)) and
xanthan (\(x\)) are both causally dependent on \(z\). The \emph{chain}
network is so-called because all the variables form a causal chain from
\(s\) to \(x\) to \(z\). (Note the directions of the arrows in the two
Bayes nets.)

Let's see what kind of data these Bayes nets produce. Let's assume that
each root node of a network (\(z\) in the common cause, \(x\) in the
chain) follows a normal distribution with mean 0 and SD 1. Each link in
a network follows a normal distribution with mean equal to the value of
its parent node and SD 1.

\begin{Shaded}
\begin{Highlighting}[]
\ImportTok{import}\NormalTok{ numpy }\ImportTok{as}\NormalTok{ np}

\NormalTok{n\_samples }\OperatorTok{=} \DecValTok{10000}

\CommentTok{\# Common cause}
\NormalTok{z\_mu }\OperatorTok{=} \DecValTok{0}
\NormalTok{sd }\OperatorTok{=} \DecValTok{1}

\NormalTok{z\_samples\_cc }\OperatorTok{=}\NormalTok{ np.zeros(n\_samples)}
\NormalTok{s\_samples\_cc }\OperatorTok{=}\NormalTok{ np.zeros(n\_samples)}
\NormalTok{x\_samples\_cc }\OperatorTok{=}\NormalTok{ np.zeros(n\_samples)}

\ControlFlowTok{for}\NormalTok{ i }\KeywordTok{in} \BuiltInTok{range}\NormalTok{(n\_samples):}
\NormalTok{  z\_samples\_cc[i] }\OperatorTok{=}\NormalTok{ np.random.normal(z\_mu, sd)}
\NormalTok{  s\_samples\_cc[i] }\OperatorTok{=}\NormalTok{ np.random.normal(z\_samples\_cc[i], sd)}
\NormalTok{  x\_samples\_cc[i] }\OperatorTok{=}\NormalTok{ np.random.normal(s\_samples\_cc[i], sd)}

\CommentTok{\# Chain}
\NormalTok{x\_mu }\OperatorTok{=} \DecValTok{0}

\NormalTok{z\_samples\_chain }\OperatorTok{=}\NormalTok{ np.zeros(n\_samples)}
\NormalTok{s\_samples\_chain }\OperatorTok{=}\NormalTok{ np.zeros(n\_samples)}
\NormalTok{x\_samples\_chain }\OperatorTok{=}\NormalTok{ np.zeros(n\_samples)}

\ControlFlowTok{for}\NormalTok{ i }\KeywordTok{in} \BuiltInTok{range}\NormalTok{(n\_samples):}
\NormalTok{  x\_samples\_chain[i] }\OperatorTok{=}\NormalTok{ np.random.normal(x\_mu, sd)}
\NormalTok{  z\_samples\_chain[i] }\OperatorTok{=}\NormalTok{ np.random.normal(x\_samples\_chain[i], sd)}
\NormalTok{  s\_samples\_chain[i] }\OperatorTok{=}\NormalTok{ np.random.normal(z\_samples\_chain[i], sd)}
\end{Highlighting}
\end{Shaded}

Because \(z\) represents an unknown variable, let's plot just \(s\) and
\(x\).

\begin{Shaded}
\begin{Highlighting}[]
\ImportTok{import}\NormalTok{ matplotlib.pyplot }\ImportTok{as}\NormalTok{ plt}
\ImportTok{import}\NormalTok{ seaborn }\ImportTok{as}\NormalTok{ sns}

\NormalTok{sns.set\_theme()}

\NormalTok{fig, axes }\OperatorTok{=}\NormalTok{ plt.subplots(}\DecValTok{1}\NormalTok{,}\DecValTok{2}\NormalTok{)}

\NormalTok{sns.scatterplot(x }\OperatorTok{=}\NormalTok{ s\_samples\_cc, y }\OperatorTok{=}\NormalTok{ x\_samples\_cc, alpha }\OperatorTok{=} \FloatTok{0.1}\NormalTok{, ax }\OperatorTok{=}\NormalTok{ axes[}\DecValTok{0}\NormalTok{])}
\NormalTok{axes[}\DecValTok{0}\NormalTok{].set\_xlabel(}\StringTok{"s"}\NormalTok{)}
\NormalTok{axes[}\DecValTok{0}\NormalTok{].set\_ylabel(}\StringTok{"x"}\NormalTok{)}
\NormalTok{axes[}\DecValTok{0}\NormalTok{].set\_title(}\StringTok{"Common cause"}\NormalTok{)}

\NormalTok{sns.scatterplot(x }\OperatorTok{=}\NormalTok{ s\_samples\_chain, y }\OperatorTok{=}\NormalTok{ x\_samples\_chain, alpha }\OperatorTok{=} \FloatTok{0.1}\NormalTok{, ax }\OperatorTok{=}\NormalTok{ axes[}\DecValTok{1}\NormalTok{])}
\NormalTok{axes[}\DecValTok{1}\NormalTok{].set\_xlabel(}\StringTok{"s"}\NormalTok{)}
\NormalTok{axes[}\DecValTok{1}\NormalTok{].set\_ylabel(}\StringTok{"x"}\NormalTok{)}
\NormalTok{axes[}\DecValTok{1}\NormalTok{].set\_title(}\StringTok{"Chain"}\NormalTok{)}

\NormalTok{plt.show()}
\end{Highlighting}
\end{Shaded}

\pandocbounded{\includegraphics[keepaspectratio]{10-causal-inference_files/figure-pdf/cell-3-output-1.pdf}}

Clearly, the data isn't identical in the two cases, but data generated
by both Bayes nets results in a strong positive correlation between
\(s\) and \(x\).

Imagine you didn't know how this data was generated and you just got one
of these plots. Could you use it to tell whether it was produced by a
common cause structure or a chain structure?

Sorry, but no. 😔 Just knowing the data are positively correlated
doesn't give you enough information to figure out how \(s\) and \(x\)
are causally related.

But what if you could manipulate the variables? That is, what if you
could \emph{intervene} on sonin levels and see what effect it had on
xanthan levels?

\begin{itemize}
\tightlist
\item
  If you increase the sonin levels 📈 and the xanthan levels also
  increase 📈, then the causal structure must be a chain.
\item
  If you increase the sonin levels 📈 and the xanthan levels don't
  change ❌, then the causal structure can't be a chain (and therefore
  must be a common cause, because that's the only other option we're
  considering).
\end{itemize}

\subsection{Graph surgery}\label{graph-surgery}

This intuition can be illustrated visually on the Bayes nets by
performing ``surgery'' on the graphs. It works like this:

\begin{enumerate}
\def\labelenumi{\arabic{enumi}.}
\tightlist
\item
  Remove all incoming connections to the variable you're intervening on.
\item
  If there's still a path between the variable you intervened on and
  another variable, then you should still expect those variables to be
  related.
\end{enumerate}

Let's apply this idea to our common cause and chain Bayes nets.

\begin{figure}[H]

{\centering \pandocbounded{\includegraphics[keepaspectratio]{images/10/graph_surgery.png}}

}

\caption{Graph surgery applied to the common cause and chain Bayes nets.
The intervention is indicated by the red arrow.}

\end{figure}%

After intervening on sonin levels (\(s\)), we remove the connection to
\(s\) in the common cause network, but no connections in the chain
network. The resulting Bayes nets show why we should expect to see a
resulting change in xanthan for the chain, but not the common cause.

\subsection{Do people intuitively understand the logic of casual
intervention?}\label{do-people-intuitively-understand-the-logic-of-casual-intervention}

\href{https://psycnet.apa.org/doi/10.1037/0278-7393.31.2.216}{A study by
Michael Waldmann and York Hagmayer} presented people with either the
common cause or the chain structure. They were told that sonin and
xanthan were hormone levels in chimps and they got some example data
that allowed them to learn that the hormone levels were positively
correlated.

Then they were either assigned to a \emph{doing} or \emph{seeing}
condition. People in the doing condition were asked to imagine that 20
chimps had their sonin levels raised (or lowered). They then predicted
how many of the chimps would have elevated xanthan levels. People in the
seeing condition got essentially the same information but just learned
that the chimps' sonin levels were high (not that it had been
intentionally raised).

Average results and model predictions are below.

\begin{figure}[H]

{\centering \pandocbounded{\includegraphics[keepaspectratio]{images/10/waldmann_hagmayer_results.png}}

}

\caption{Results from Waldmann \& Hagmayer (2005), Experiment 2.}

\end{figure}%

People's judgments mostly followed those of the Bayes net model
predictions. In the common cause case, when the sonin levels are
artificially raised, the relationship between sonin and xanthan is
decoupled, so the model reverts to a base rate prediction about xanthan
levels (and so did people, for the most part). But when just observing
elevated sonin levels, the model should expect the positive relationship
to hold.

In the chain case, the predictions are the same for seeing or doing,
because intervening doesn't change anything about the relationship
between sonin and xanthan. People's judgments indicate that they
understood this.

\section{Causal structure and strength 🏗💪}\label{structure-strength}

Bayes nets can also account for how people judge the strength of
evidence for a causal relationship after seeing some data. This was the
idea that Tom Griffiths and Josh Tenenbaum explored in
\href{https://cocosci.princeton.edu/tom/papers/causal4.pdf}{a 2005
computational study}.

Here's the basic problem they considered. Suppose researchers perform an
experiment with rats to test whether a drug causes a gene to be
expressed. A control group of rats doesn't get the injection and the
experimental group does. They record the number of rats in each group
that express the gene.

Here are some possible results from an experiment with 8 rats in each
group. The table shows how many rats in each group expressed the gene.

\begin{longtable}[]{@{}ll@{}}
\toprule\noalign{}
Control 🐀 & Experimental 🐀 \\
\midrule\noalign{}
\endhead
\bottomrule\noalign{}
\endlastfoot
6/8 🧬 & 8/8 🧬 \\
4/8 🧬 & 6/8 🧬 \\
2/8 🧬 & 4/8 🧬 \\
0/8 🧬 & 2/8 🧬 \\
\end{longtable}

In each of these hypothetical experiments, how strongly would you say
that the drug causes the gene to be expressed?

These are a few of the cases that were included in an experiment
conducted by Marc Buehner and Patricia Cheng. Here's the full set of
averaged human results.

\begin{figure}[H]

{\centering \pandocbounded{\includegraphics[keepaspectratio]{images/10/buehner_cheng_results.png}}

}

\caption{Results from Buehner \& Cheng (1997), Experiment 1B, reproduced
by Griffiths \& Tenenbaum (2005). p(e+\textbar c+) is the probability of
the effect being present (e.g., the gene being expressed) given that the
cause is present (e.g., the drug was administered). p(e+\textbar c-)
means the probability of the effect being present given that the cause
is absent.}

\end{figure}%

Focusing just on the cases in the table, on average, people judged that
the drug was less likely to have a causal effect on the gene as the
total number of rats expressing the gene decreased, even when the
difference in number of rats expressing the gene between conditions was
held constant.

\subsection{The causal support model}\label{the-causal-support-model}

Maybe people reason about these problems by performing a kind of model
selection between the two Bayes nets below.

\begin{figure}[H]

{\centering \includegraphics[width=0.75\linewidth,height=\textheight,keepaspectratio]{images/10/griffiths_tenenbaum_graphs.png}

}

\caption{Causal inference as model selection. In one model, there is a
connection from the potential cause to the effect; in the other, the two
variables are unconnected. Image from Griffiths \& Tenenbaum (2005).}

\end{figure}%

These Bayes nets each have three variables: an effect \(E\), a cause
\(C\), and a background cause \(B\). For our problem, the effect and
cause refer to the gene and the drug. The inclusion of the background
cause is to account for unknown factors that might cause the gene to be
expressed without the drug.

The problem people are faced with is deciding which of these two models
is best supported by the data \(D\) -- the number of times the effect
occurred with and without the potential cause.

This can be done with Bayesian inference:

\[
P(\text{Graph } i) \propto P(D|\text{Graph } i) P(\text{Graph } i)
\]

Because there are only two possible networks, we can compute the
relative evidence for one Bayes net over the other as a ratio. We can
then take the log of the expression to simplify it:

\[
\log \frac{P(\text{Graph } 1|D)}{P(\text{Graph } 0|D)} = \log \frac{P(D | \text{Graph } 1)}{P(D | \text{Graph } 0)} + \log \frac{P(\text{Graph } 1)}{P(\text{Graph } 0)}
\]

Regardless of what prior probabilities we assign to the two graphs, the
relative evidence for one graph over the other is entirely determined by
the log-likelihood ratio. This is defined as \textbf{causal support}:
\(\log \frac{P(D | \text{Graph } 1)}{P(D | \text{Graph } 0)}\).

Computing \(P(D | \text{Graph } 1)\) requires fully specifying the Bayes
net. We'll assume that \(P(E|B) = w_0\) and \(P(E|C) = w_1\). When both
\(B\) and \(C\) are present, we'll assume they contribute independently
to causing \(E\), and therefore operate like a probabilistic OR
function:

\[
P(e^+|b,c; w_o, w_1) = 1 - (1-w_0)^b (1-w_1)^c
\]

Here, when the B or C are present, \(b\) and \(c\) are set to 1, and
when they are absent, \(b\) and \(c\) are set to 0.

The likelihood for Graph 1 is therefore

\[
P(D | w_0, w_1, \text{Graph } 1) = \prod_{e,c} P(e|b^+,c; w_o, w_1)^{N(e,c)}
\]

where the product is over the possible settings of \(e\) and \(c\)
(effect absent/cause absent, effect absent/cause present, \ldots) and
the \(N(e,c)\) values are counts of times that these outcomes happened
in the data \(D\).

Here's a function to compute this.

\begin{Shaded}
\begin{Highlighting}[]
\KeywordTok{def}\NormalTok{ compute\_likelihood(data, w0, w1, graph):}
  \CommentTok{\textquotesingle{}\textquotesingle{}\textquotesingle{}}
\CommentTok{  Computes likelihood of data for a given graph}
\CommentTok{     }
\CommentTok{  Args:}
\CommentTok{    data (list): observation counts in this order:}
\CommentTok{      N(c{-},e{-}), N(c{-},e+), N(c+,e{-}), N(c+,e+)}
\CommentTok{    w0 (float): probability of background cause producing effect}
\CommentTok{    w1 (float): probability of cause of interest producing effect}
\CommentTok{      graph (int): 0 (Graph 0) or 1 (Graph 1)}

\CommentTok{    Returns:}
\CommentTok{      (float): probability of data }
\CommentTok{  \textquotesingle{}\textquotesingle{}\textquotesingle{}}
  
  \ControlFlowTok{if}\NormalTok{ graph }\OperatorTok{==} \DecValTok{0}\NormalTok{:}
    \CommentTok{\# e{-}}
\NormalTok{    p }\OperatorTok{=}\NormalTok{ (}\DecValTok{1}\OperatorTok{{-}}\NormalTok{w0)}\OperatorTok{**}\NormalTok{(data[}\DecValTok{0}\NormalTok{]}\OperatorTok{+}\NormalTok{data[}\DecValTok{2}\NormalTok{])}
    \CommentTok{\# e+}
\NormalTok{    p }\OperatorTok{=}\NormalTok{ p }\OperatorTok{*}\NormalTok{ w0}\OperatorTok{**}\NormalTok{(data[}\DecValTok{1}\NormalTok{]}\OperatorTok{+}\NormalTok{data[}\DecValTok{3}\NormalTok{])}
  \ControlFlowTok{elif}\NormalTok{ graph }\OperatorTok{==} \DecValTok{1}\NormalTok{:}
    \CommentTok{\# c{-}, e{-}}
\NormalTok{    p }\OperatorTok{=}\NormalTok{ (}\DecValTok{1}\OperatorTok{{-}}\NormalTok{w0)}\OperatorTok{**}\NormalTok{data[}\DecValTok{0}\NormalTok{]}
    \CommentTok{\# c{-}, e+}
\NormalTok{    p }\OperatorTok{=}\NormalTok{ p }\OperatorTok{*}\NormalTok{ w0}\OperatorTok{**}\NormalTok{data[}\DecValTok{1}\NormalTok{]}
    \CommentTok{\# c+, e{-}}
\NormalTok{    p }\OperatorTok{=}\NormalTok{ p }\OperatorTok{*}\NormalTok{ (}\DecValTok{1} \OperatorTok{{-}}\NormalTok{ (w0 }\OperatorTok{+}\NormalTok{ w1 }\OperatorTok{{-}}\NormalTok{ w0}\OperatorTok{*}\NormalTok{w1))}\OperatorTok{**}\NormalTok{data[}\DecValTok{2}\NormalTok{]}
    \CommentTok{\# c+, e+}
\NormalTok{    p }\OperatorTok{=}\NormalTok{ p }\OperatorTok{*}\NormalTok{ (w0 }\OperatorTok{+}\NormalTok{ w1 }\OperatorTok{{-}}\NormalTok{ w0}\OperatorTok{*}\NormalTok{w1)}\OperatorTok{**}\NormalTok{data[}\DecValTok{3}\NormalTok{]}
  \ControlFlowTok{else}\NormalTok{:}
    \CommentTok{\# error!}
    \ControlFlowTok{return}\NormalTok{(}\DecValTok{0}\NormalTok{)}
  
  \ControlFlowTok{return}\NormalTok{(p)}
\end{Highlighting}
\end{Shaded}

Causal support doesn't actually depend directly on the parameters
\(w_0\) and \(w_1\). The reason is that we ultimately don't care what
the values of these parameters are because we just want to draw an
inference at a higher level about the best-fitting Bayes net.

Mathematically speaking, \(w_0\) and \(w_1\) are averaged out of the
model, an idea we first saw in \hyperref[generalizing]{Chapter 3}. We
can accomplish this using Monte Carlo approximation, introduced in
\hyperref[monte-carlo]{Chapter 8}.

\begin{Shaded}
\begin{Highlighting}[]
\KeywordTok{def}\NormalTok{ estimate\_likelihood(graph\_number, data): }
  \CommentTok{\textquotesingle{}\textquotesingle{}\textquotesingle{}}
\CommentTok{  Returns an estimate of the probability of observing the data}
\CommentTok{    under the specified graph using Monte Carlo estimation.}
\CommentTok{     }
\CommentTok{  Args:}
\CommentTok{    graph\_number (int): either 0 (Graph 0) or 1 (Graph 1)}
\CommentTok{    data (list): observation counts in this order:}
\CommentTok{      N(c{-},e{-}), N(c{-},e+), N(c+,e{-}), N(c+,e+)}

\CommentTok{  Returns:}
\CommentTok{    (float): probability of observing data}
\CommentTok{  \textquotesingle{}\textquotesingle{}\textquotesingle{}}
  \ImportTok{from}\NormalTok{ numpy.random }\ImportTok{import}\NormalTok{ default\_rng}
  
\NormalTok{  n\_samples }\OperatorTok{=} \DecValTok{5000}
\NormalTok{  rng }\OperatorTok{=}\NormalTok{ default\_rng(}\DecValTok{2022}\NormalTok{)}
\NormalTok{  mc\_samples }\OperatorTok{=}\NormalTok{ np.zeros(n\_samples)}
  
  \ControlFlowTok{for}\NormalTok{ i }\KeywordTok{in} \BuiltInTok{range}\NormalTok{(n\_samples):}
    \CommentTok{\# Sample values for w0 and w1 from a uniform distribution}
\NormalTok{    w0 }\OperatorTok{=}\NormalTok{ rng.random()}
\NormalTok{    w1 }\OperatorTok{=}\NormalTok{ rng.random()}
    
\NormalTok{    mc\_samples[i] }\OperatorTok{=}\NormalTok{ compute\_likelihood(data, w0, w1, graph\_number)}
  
  \ControlFlowTok{return}\NormalTok{(}\DecValTok{1}\OperatorTok{/}\NormalTok{n\_samples }\OperatorTok{*}\NormalTok{ np.}\BuiltInTok{sum}\NormalTok{(mc\_samples))}
\end{Highlighting}
\end{Shaded}

Finally, let's put it all together by writing a function that computes
causal support.

\begin{Shaded}
\begin{Highlighting}[]
\KeywordTok{def}\NormalTok{ causal\_support(data):}
  \CommentTok{\textquotesingle{}\textquotesingle{}\textquotesingle{}}
\CommentTok{  Returns a causal support value for a given data set.}
\CommentTok{  Causal support is a measure of how strongly a data}
\CommentTok{  set indicates that there is evidence for a causal}
\CommentTok{  effect.}
\CommentTok{     }
\CommentTok{  Parameters:}
\CommentTok{   data (list): observation counts in this order:}
\CommentTok{    N(c{-},e{-}), N(c{-},e+), N(c+,e{-}), N(c+,e+)}

\CommentTok{  Returns:}
\CommentTok{    (float): strength of evidence for causal relationship}
\CommentTok{  \textquotesingle{}\textquotesingle{}\textquotesingle{}}
  \ImportTok{import}\NormalTok{ numpy }\ImportTok{as}\NormalTok{ np}
  
  \ControlFlowTok{return}\NormalTok{ (np.log(estimate\_likelihood(}\DecValTok{1}\NormalTok{,data}\OperatorTok{=}\NormalTok{data) }\OperatorTok{/} 
\NormalTok{                 estimate\_likelihood(}\DecValTok{0}\NormalTok{,data}\OperatorTok{=}\NormalTok{data)))}
\end{Highlighting}
\end{Shaded}

Now let's see what the model predicts for the
\hyperref[structure-strength]{four cases we considered earlier}.

\begin{Shaded}
\begin{Highlighting}[]
\NormalTok{causal\_support\_predictions }\OperatorTok{=}\NormalTok{ [causal\_support([}\DecValTok{2}\NormalTok{,}\DecValTok{6}\NormalTok{,}\DecValTok{0}\NormalTok{,}\DecValTok{8}\NormalTok{]),}
\NormalTok{                              causal\_support([}\DecValTok{4}\NormalTok{,}\DecValTok{4}\NormalTok{,}\DecValTok{2}\NormalTok{,}\DecValTok{6}\NormalTok{]),}
\NormalTok{                              causal\_support([}\DecValTok{6}\NormalTok{,}\DecValTok{2}\NormalTok{,}\DecValTok{4}\NormalTok{,}\DecValTok{4}\NormalTok{]),}
\NormalTok{                              causal\_support([}\DecValTok{8}\NormalTok{,}\DecValTok{0}\NormalTok{,}\DecValTok{6}\NormalTok{,}\DecValTok{2}\NormalTok{])]}
                              
\NormalTok{labels }\OperatorTok{=}\NormalTok{ [}\StringTok{"6/8, 8/8"}\NormalTok{, }\StringTok{"4/8, 6/8"}\NormalTok{, }\StringTok{"2/8, 4/8"}\NormalTok{, }\StringTok{"0/8, 2/8"}\NormalTok{]}
\NormalTok{fig, ax }\OperatorTok{=}\NormalTok{ plt.subplots()}
\NormalTok{sns.barplot(x}\OperatorTok{=}\NormalTok{labels, y}\OperatorTok{=}\NormalTok{causal\_support\_predictions)}
\NormalTok{ax.set\_xlabel(}\StringTok{"Data"}\NormalTok{)}
\NormalTok{ax.set\_ylabel(}\StringTok{"Model prediction"}\NormalTok{)}

\NormalTok{plt.show()}
\end{Highlighting}
\end{Shaded}

\pandocbounded{\includegraphics[keepaspectratio]{10-causal-inference_files/figure-pdf/cell-7-output-1.pdf}}

The labels on the x-axis indicate the control condition counts, followed
by the experimental condition counts.

You can see that the model's predictions largely follow the
\hyperref[structure-strength]{pattern in the data from earlier}. The
model favors Graph 1 in the two leftmost cases, is essentially uncertain
in the third case, and begins to think the evidence favors no causal
relationship in the rightmost case.

The rest of the Griffiths \& Tenenbaum paper shows how causal support is
able to capture some subtle aspects of people's causal judgments that
other models that don't incorporate \emph{both} structure and strength
fail to capture.




\end{document}
